\documentclass[11pt,a4paper]{article}
\usepackage[margin=27mm]{geometry}
\usepackage[T1]{fontenc}
\usepackage{lmodern,microtype,amsmath,amssymb,amsthm,mathtools}
\usepackage{enumitem,booktabs,array,longtable,needspace}
\usepackage[hidelinks]{hyperref}
\usepackage[nameinlink,noabbrev]{cleveref}
\setlength{\parindent}{1.2em}
\setlength{\parskip}{2pt}
\setlist[enumerate]{label=\textup{(\roman*)},leftmargin=2em}
\newtheorem{theorem}{Theorem}[section]
\newtheorem{proposition}[theorem]{Proposition}
\newtheorem{lemma}[theorem]{Lemma}
\newtheorem{corollary}[theorem]{Corollary}
\theoremstyle{definition}
\newtheorem{definition}[theorem]{Definition}
\theoremstyle{remark}
\newtheorem{remark}[theorem]{Remark}
\title{Positive invariance entropy under\\uniform state convergence}
\author{}
\date{}
\begin{document}
\maketitle
\begin{abstract}
We construct smooth control systems with globally invertible controlled maps and full-dimensional compact convex constraints for which exact, outer, and reference-trajectory tracking invariance entropies are equal and strictly positive, although all feasible state trajectories converge uniformly to one common fixed state. For a general affine system viable in a convex body, adjoining a decreasing radial coordinate gives exact finite-horizon transference to a solid cone. On every compact initial set of positive ambient volume, the outer and tracking rates are the positive spectral sum of \(\rho A\), where \(\rho\) is the radial derivative at the apex; a uniform interior viability condition gives the exact rate as the positive spectral sum of \(A\). In dimension two these formulas persist under an infinite-dimensional class of smooth nonlinear perturbations preserving the first jet. The proof combines uniform small-radius control blocks, sectional-volume estimates, and a reduction of tracking to a shorter exact-control horizon. Choosing the radial map \(\operatorname{arsinh}\) gives equal positive rates despite subexponential convergence of every feasible trajectory. A common Lyapunov condition identifies the resulting entropy separation: the state-trajectory shift and the input fibers have zero entropy, all invariant lift measures are supported on the common fixed state, and the total lift entropy comes entirely from the input shift. The construction concerns forward controlled invariant constraints, not recurrent control sets.
\end{abstract}
\noindent\textbf{2020 Mathematics Subject Classification.} Primary 37B40, 93C55; Secondary 37B05.\\
\textbf{Keywords.} Invariance entropy; outer invariance entropy; Li--Yorke pairs; cone constraints; nonlinear perturbations; control skew product; relative entropy.

\section{Introduction}

Invariance entropy measures the number of control signals that must be available in advance to maintain a prescribed state constraint. Positive topological entropy, by contrast, measures the orbit complexity of one compact dynamical system and forces Li--Yorke chaos for continuous surjections \cite{BGKM}. The optimizations are different: the signal assigned to an initial state may need arbitrarily fine precision even when every feasible state trajectory approaches the same point.

We study this distinction on full-dimensional compact convex constraints for smooth invertible controlled maps. Our principal examples have equal, finite, strictly positive exact, outer, and tracking invariance entropies, while every pair of feasible state trajectories is uniformly asymptotic, even when the two trajectories use different inputs. We prove a spectral formula in arbitrary dimension and show that its planar version persists under an infinite-dimensional class of nonlinear perturbations preserving the first jet at the common fixed state.

\subsection{Main results}

Here \(h_{\mathrm{inv}}\) denotes exact finite-catalogue invariance entropy, \(h_{\mathrm{out}}\) uses ambient neighborhoods of the constraint, and \(h_{\mathrm{inv}}^+\) uses finite catalogues of viable reference-input pairs. Their precise quantifiers and limit orders are given in \cref{sec:definitions}. All logarithms are natural.

Let \(A\in GL(d,\mathbb R)\), let \(B\subset\mathbb R^d\) be a compact convex body with \(0\in\operatorname{int}B\), and let \(U\subset\mathbb R^d\) be nonempty and compact. Assume that
\(\forall y\in B\ \exists v\in U:\ Ay+v\in B\).
Let \(R:\mathbb R\to\mathbb R\) be an increasing smooth diffeomorphism with \(R(0)=0\) and \(0<R(s)<s\) for \(0<s\leq1\). Write \(\rho=R'(0)\in(0,1]\), put
\(Q_B=\{(s,sy):0\leq s\leq1,\ y\in B\}\), and consider
\begin{align}\label{eq:system-intro}
 f((s,z),v)=\left(R(s),\frac{R(s)}s(Az+sv)\right),
\end{align}
using the smooth extension \(R(s)/s=\rho\) at \(s=0\). For \(L\in GL(d,\mathbb R)\), write
\begin{align}\label{eq:chi}
 \chi^+(L)=\sum_{j=1}^d\max\{0,\log|\lambda_j(L)|\},
\end{align}
counting algebraic multiplicities. In the next theorem, superscripts \(f\) and \(g\) distinguish \eqref{eq:system-intro} from its normalized affine system \(g(y,v)=Ay+v\).

\begin{theorem}[Cone transference and spectral formulas]\label{thm:cone}
Under the preceding assumptions:
\begin{enumerate}
\item \(Q_B\) is a full-dimensional compact convex controlled invariant set, and every \(f_v\) is a global \(C^\infty\) diffeomorphism of \(\mathbb R^{d+1}\). For every integer \(n\geq0\), including possible infinite values,
\begin{align}\label{eq:transference}
 r_{\mathrm{inv}}^{f}(n,Q_B,Q_B)=r_{\mathrm{inv}}^{g}(n,B,B).
\end{align}
\item For every compact \(K\subset Q_B\) with \(\operatorname{vol}_{d+1}(K)>0\),
\begin{align}\label{eq:spectral}
 h_{\mathrm{out}}^{f}(K,Q_B)=h_{\mathrm{inv}}^{f,+}(K,Q_B)=\chi^+(\rho A),
 \qquad
 h_{\mathrm{inv}}^{f}(K,Q_B)\geq\chi^+(A).
\end{align}
For each fixed \(\varepsilon>0\), the normalized logarithms defining the outer and tracking counts have actual limits, both equal to \(\chi^+(\rho A)\).
\item Suppose, additionally, that there is \(\eta_0>0\) such that
\begin{align}\label{eq:margin}
 \forall y\in B\quad\exists v\in U:\quad
 \operatorname{dist}(Ay+v,\mathbb R^d\setminus B)\geq\eta_0.
\end{align}
Then every such \(K\) satisfies
\begin{align}\label{eq:exact-spectral}
 \lim_{n\to\infty}\frac1n\log r_{\mathrm{inv}}^{f}(n,K,Q_B)
       =h_{\mathrm{inv}}^{f}(K,Q_B)=\chi^+(A).
\end{align}
\end{enumerate}
\end{theorem}

\begin{theorem}[Nonlinear perturbations preserving the entropy rates]\label{thm:nonlinear}
Let \(q>1\), let \(U\subset\mathbb R\) be finite and nonempty, and suppose that for some \(\eta>0\)
\begin{align}\label{eq:nonlinear-margin}
 \forall y\in[-1,1]\quad\exists v\in U:
       |qy+v|\leq1-2\eta.
\end{align}
Let \(R\) satisfy the radial assumptions above, put \(\rho=R'(0)\), and extend \(t(s)=R(s)/s\) smoothly at zero. For every \(v\in U\), let \(P_v:\mathbb R^2\to\mathbb R\) be smooth, with
\begin{align}\label{eq:zero-jet}
 P_v(0,0)=0,\qquad DP_v(0,0)=0,\qquad
 \|P\|_{C_b^2}<\min\{\eta,q/2\},
\end{align}
where \(\|P\|_{C_b^2}\) is the maximum, over \(v\in U\), of the global suprema of \(|P_v|\), the Euclidean operator norm of \(DP_v\), and that of \(D^2P_v\). Define
\begin{align}\label{eq:nonlinear-system}
 f_v(s,z)=\bigl(R(s),t(s)[qz+sv+P_v(s,z)]\bigr),
 \qquad Q=\{(s,z):0\leq s\leq1,\ |z|\leq s\}.
\end{align}
Then each \(f_v\) is a global smooth diffeomorphism, \(Q\) is controlled invariant, and for every compact \(K\subset Q\) of positive area,
\begin{align}\label{eq:nonlinear-entropies}
 \lim_{n\to\infty}\frac1n\log r_{\mathrm{inv}}(n,K,Q)
 &=h_{\mathrm{inv}}(K,Q)=\log q,\\
 h_{\mathrm{out}}(K,Q)=h_{\mathrm{inv}}^+(K,Q)
 &=\max\{0,\log(\rho q)\}.
\end{align}
All feasible states at time \(n\), uniformly over initial states and inputs, lie in \(R^n(1)Q\). Thus every pair of feasible state trajectories is asymptotic.
\end{theorem}

The perturbations form a nonempty open ball relative to the space of bounded \(C^2\) perturbations with zero first jet, intersected with the smooth functions. The radial map and the control alphabet are held fixed. This is a persistence statement in a specified infinite-dimensional class; it does not assert openness under arbitrary changes of the radial dynamics or of the linearization.
\begin{corollary}[Equal positive control entropies with zero state entropy]\label{cor:joint}
Set \(R(s)=\operatorname{arsinh}s\), \(q=2\), \(U=\{-3/2,0,3/2\}\), and
\begin{align}\label{eq:explicit-nonlinear}
 P_v(s,z)=e_v\tanh(s)\tanh(z),\qquad
 \max_{v\in U}|e_v|<1/64.
\end{align}
The maps \eqref{eq:nonlinear-system} are global smooth diffeomorphisms, and for every compact \(K\subset Q\) of positive area,
\begin{align}
 h_{\mathrm{inv}}(K,Q)=h_{\mathrm{out}}(K,Q)
 =h_{\mathrm{inv}}^+(K,Q)=\log2.
\end{align}
Every pair of feasible state trajectories has distance at most
\(2R^n(1)\to0\) at time \(n\), with \(R^n(1)\sim\sqrt{3/n}\).
The true constrained lift has total topological entropy \(\log3\), uniform entropy over input fibers zero, and state-trajectory entropy zero. Every invariant lift measure is concentrated on the common fixed state and has zero conditional Kolmogorov--Sinai entropy over the input factor.
\end{corollary}

The two-symbol unperturbed example and its exact finite-horizon counts appear in \cref{thm:digits,cor:equal}; the proof of \cref{cor:joint} is given in \cref{sec:joint-proof}. The equalities in this corollary show that the phenomenon persists when all three control entropies are positive. It is not explained by an exact-versus-outer entropy gap.

The proofs have three separate components. Linear Bowen coverings and unstable sectional volumes yield \cref{thm:cone}, without diagonalizability or a hyperbolicity assumption. A finite-horizon reduction for tracking catalogues, proved in \cref{prop:shortening}, supplies the nonlinear upper bound in \cref{thm:nonlinear}; uniform small-radius blocks give its exact upper bound. Finally, \cref{thm:lyapunov} uses a common Lyapunov inequality to identify the state factor, the input fibers, all invariant measures, and the eventual core of the true constrained lift. Positivity of control entropy is supplied by the spectral and perturbation theorems, not by the Lyapunov inequality alone. \Cref{cor:rates} also realizes every prescribed pair of finite rates \(a>0\), \(0\leq b\leq a\), as exact and outer invariance entropies.

\subsection{Relation to earlier work}

Invariance entropy was introduced and developed in \cite{CK2009,K2009,K2011}. Exact, strict, outer, and strong terminology varies across sources; the tracking quantity used here is the discrete-time counterpart of the reference-trajectory definition in \cite[Definition 3.1.1]{K2009}. Spectral expressions, linear coverings, and unstable-volume lower estimates are established tools, including the linear and inhomogeneous bilinear calculations in \cite[Theorem 4.1.2, Proposition 4.1.9, and Theorem 4.1.10]{K2009} and the strict lower-bound theory \cite{KStrict2011}. We claim neither a new volume method nor novelty for the logarithm of an expanding determinant.

The upper bounds for control sets and the hyperbolic formula of Da Silva and Kawan \cite{DSK2016} concern complete approximate controllability and full-time invariant lifts. Our entire solid cone is only forward controlled invariant: a positive radial coordinate decreases under every input, and the full-time feasible state core reduces to the apex. Thus replacing the cone by its core loses the positive-volume initial sets used here. The equality examples also have a neutral radial derivative at the apex, outside a uniformly hyperbolic splitting without center. These are substantive differences in hypotheses, beyond the discrete-versus-continuous time distinction.

The broader separation between control information and dynamical complexity already appears in the partially hyperbolic pressure bound of Kawan and Da Silva \cite{KDS2018}, and in Kawan's discrete hyperbolic theory \cite{KChaos}, where unstable volume growth and measure-theoretic fiber entropy enter as distinct terms. Those works are conceptual precedents. Their full-time isolated lifts and hyperbolic or partially hyperbolic assumptions do not directly provide the exact/outer/tracking calculation on the forward cone considered here. In particular, the present joint phenomenon should not be read as the first observation that invariance entropy differs from state entropy.

Yao's preprint \cite[Theorem 3.1]{Yao2026} gives positive exact invariance entropy and zero approximate invariance entropy on a thin constraint. Its equations already imply that feasible pair distances have limits: two state coordinates decay and a third is constant. We therefore do not claim as new the weak statement that positive exact invariance entropy can coexist with the absence of state Li--Yorke pairs. Our focus is the simultaneous positive outer and tracking rates, the solid convex constraint, and the explicit nonlinear class.

The formulas are also related to stabilization and network data rates \cite{CH2021,KD2016}. Those problems optimize controllers for prescribed stabilization or synchronization objectives. Here feasibility itself forces convergence for every viable input; no selected controller is required to achieve an externally prescribed decay rate. Affine homogenization and projective representations are familiar \cite{CST2023}. However, the coordinate \(y=z/s\) deletes the apex and is singular there, so compact-constraint entropy cannot simply be transferred through that coordinate change. Nor are distinct common-fixed-point maps translations of a single Lie-group automorphism; \cref{rem:lie} explains the obstruction to directly applying \cite{CCS2021}.

Equi-invariability and invariance complexity study nearby initial conditions, finite control families, and possible departure from a constraint \cite{ZCH2021,ZC2023,CZ2024}. Departure of a nearby trajectory is compatible with convergence of all trajectories that remain feasible. The measurable data-rate and entropy-dimension theories \cite{WH2024,ZH2025,CHZ2026} concern further optimizations and growth scales, rather than the conditional Kolmogorov--Sinai entropy of the lift computed below. No theorem-level exclusion of their other results is inferred merely from those descriptions.

Zhang's relative theory \cite[Theorems 3.9, 3.10, and 4.2]{Zhang2006} obtains asymptotic pairs and fiber Li--Yorke chaos from positive conditional entropy. Its uniform fiber formula \cite[Lemma 4.3]{Zhang2006} has the same order of supremum and limits as \eqref{eq:fiber-definition}, although its ambient systems are homeomorphisms. Our fiber rate is zero, and our true forward lift need not be surjective; positive control catalogue growth supplies neither missing hypothesis. Random Li--Yorke theory \cite{LTZ2024} likewise requires its specified state entropy. Invariance pressure and partition formulas require their stated realizability and regularity assumptions \cite{YCYZ2025,ZHZ2023}; \cref{app:partition} explains the relevant interfaces. Recurrence entropy \cite{SM2026} and general nonautonomous skew-product constructions \cite{CF2026} provide additional context without identifying control entropy with the state-trajectory factor.

The boundary of the results is precise. The constraints constructed here are not approximately controllable control sets. The persistence theorem fixes the radial dynamics and the first jet; it is not a claim of unrestricted structural stability. We prove no negative conclusion for recurrent control sets, and no general statement about chaotic dynamics outside the feasible constraint. The main assertion is that positive information cost for maintaining viability can coexist with uniformly vanishing state separation in the explicitly described smooth systems.

\section{Definitions and elementary estimates}\label{sec:definitions}

\subsection{Control catalogues and state chaos}

Let \(X\) be a metric space and let \(U\) be a nonempty compact metric space. Throughout the main text \(f:X\times U\to X\) is continuous, \(\mathcal U=U^{\mathbb N_0}\) has its product topology, and every sequence is admissible. The uniquely defined forward solution is
\begin{align*}
 \varphi(0,x,u)=x,\qquad
 \varphi(n+1,x,u)=f(\varphi(n,x,u),u_n).
\end{align*}
No negative-time admissibility is presumed, even when individual maps \(f_v=f(\,\cdot\,,v)\) are invertible.
A nonempty compact set \(Q\subset X\) is \emph{controlled invariant} if
\begin{align*}
 \forall x\in Q\quad\exists u\in\mathcal U\quad
 \forall n\in\mathbb N_0:\ \varphi(n,x,u)\in Q.
\end{align*}
With unrestricted concatenation this is equivalent to the one-step condition
\(\forall x\in Q\ \exists v\in U:\ f(x,v)\in Q\).
We never use the term strong invariance to mean an unspecified strengthening of this statement. The different universal condition
\(\forall x\in Q\ \forall v\in U:\ f(x,v)\in Q\)
is not assumed.

For a nonempty compact \(K\subset Q\), an \((n,K,Q)\)-spanning family is a finite set \(\mathcal S\subset\mathcal U\) satisfying
\begin{align}\label{eq:quantifiers}
 \forall x\in K\quad\exists u\in\mathcal S\quad
 \forall k\in\{0,\ldots,n\}:\ \varphi(k,x,u)\in Q.
\end{align}
The family is fixed before the initial state is selected; only its assigned member may depend on that state. Its minimum cardinality is \(r_{\mathrm{inv}}(n,K,Q)\), with value \(+\infty\) if no finite family exists. Only the first \(n\) symbols matter. All logarithms are natural, and \(\log(+\infty)=+\infty\). Define
\begin{align}\label{eq:exact}
 h_{\mathrm{inv}}(K,Q)=\limsup_{n\to\infty}\frac1n\log r_{\mathrm{inv}}(n,K,Q).
\end{align}
We call this \emph{exact} invariance entropy; it is also called strict invariance entropy in some continuous-time sources. It does not require interior containment.

For \(Q^\varepsilon=\{x\in X:\operatorname{dist}(x,Q)<\varepsilon\}\), define
\begin{align}\label{eq:outer}
 h_{\mathrm{out}}(K,Q)=\lim_{\varepsilon\downarrow0}\limsup_{n\to\infty}
 \frac1n\log r_{\mathrm{inv}}(n,K,Q^\varepsilon).
\end{align}
The outer limit exists in \([0,+\infty]\) by monotonicity. The time limit is taken first. Neighborhoods are measured in the ambient state metric, not in a rescaled projective metric. In particular \(h_{\mathrm{out}}(K,Q)\leq h_{\mathrm{inv}}(K,Q)\).

We also use the discrete-time counterpart of strong invariance entropy in \cite[Definition 3.1.1]{K2009}. A \emph{tracking catalogue} consists of finitely many pairs \((a_i,u^i)\), \(a_i\in Q\), whose reference trajectories are feasible in \(Q\) for all nonnegative times, and for which
\begin{align*}
 \forall x\in K\quad\exists i\quad
 \max_{0\leq k\leq n}
 d(\varphi(k,x,u^i),\varphi(k,a_i,u^i))<\varepsilon.
\end{align*}
Let \(r_{\mathrm{inv}}^{+}(n,\varepsilon,K,Q)\) be the least number of such pairs and put
\begin{align*}
 h_{\mathrm{inv}}^{+}(K,Q)=\lim_{\varepsilon\downarrow0}\limsup_{n\to\infty}
 \frac1n\log r_{\mathrm{inv}}^{+}(n,\varepsilon,K,Q).
\end{align*}
Requiring reference viability only up to \(n\) gives the same number: a feasible finite reference trajectory can be continued from its endpoint by controlled invariance. The \emph{actual} trajectory in this definition is required to track its reference, and need not itself stay in \(Q\). Forgetting reference states gives
\begin{align}\label{eq:outer-tracking}
 r_{\mathrm{inv}}(n,K,Q^\varepsilon)\leq r_{\mathrm{inv}}^{+}(n,\varepsilon,K,Q).
\end{align}
There is no general identification of tracking entropy with exact entropy in our definitions. Measure-theoretic invariance entropy, including partition and quasi-stationary versions \cite{WHS2019,NWH2022,CQuasi2018}, is another object. We do not calculate it by computing a conditional Kolmogorov--Sinai entropy.

A pair of feasible state trajectories \(x_n,y_n\) is a \emph{state Li--Yorke pair} if
\begin{align*}
 \liminf_{n\to\infty}d(x_n,y_n)=0,\qquad
 \limsup_{n\to\infty}d(x_n,y_n)>0.
\end{align*}
The same input may be imposed on both trajectories, or distinct inputs may be allowed. For a scrambled set with state-dependent open-loop inputs, a single assignment \(x\mapsto u_x\) must first be fixed on the whole set; it cannot be reselected for each pair. A unified causal controller, including one with memory, gives another fixed rule. Our negative conclusion quantifies over \emph{all} pairs of feasible trajectories with arbitrary separate inputs, and hence covers all these more restrictive choices.

\subsection{The true lift and its entropies}

Let \(\sigma:\mathcal U\to\mathcal U\) be the one-sided shift and define
\begin{align}\label{eq:lift}
 \Lambda_Q=\{(u,x)\in\mathcal U\times Q:
       \varphi(n,x,u)\in Q\text{ for all }n\geq0\},\quad
 F(u,x)=(\sigma u,f(x,u_0)).
\end{align}
Each finite-time solution is continuous in \((u,x)\). Therefore \(\Lambda_Q\) is closed in the compact space \(\mathcal U\times Q\), and \(F\) is a continuous self-map of \(\Lambda_Q\). It is positively invariant; surjectivity does not follow. The input projection \(\pi(u,x)=u\) satisfies \(\pi F=\sigma\pi\). Its image need not in general be all of \(\mathcal U\).

Use a bounded metric \(d_U\leq1\) on \(U\), and the compatible input metric
\begin{align*}
 d_{\mathcal U}(u,v)=\sum_{j=0}^\infty 2^{-j-1}d_U(u_j,v_j).
\end{align*}
We equip the lift with the maximum of the input metric and the bounded state metric \(\min\{1,d\}\). For a compact system \(T:Z\to Z\), let \(s_n(T,\varepsilon)\) denote the maximum cardinality of a set separated by more than \(\varepsilon\) in the Bowen metric
\(\max_{0\leq k<n}d_Z(T^kz,T^kw)\). Then
\begin{align*}
 h_{\mathrm{top}}(T)=\lim_{\varepsilon\downarrow0}\limsup_{n\to\infty}
          n^{-1}\log s_n(T,\varepsilon).
\end{align*}
For nonempty input fibers use the same Bowen metric restricted to \(\pi^{-1}(u)\), and define the uniform fiber growth rate
\begin{align}\label{eq:fiber-definition}
 h_{\mathrm{fib}}(F\mid\pi)=
 \lim_{\varepsilon\downarrow0}\limsup_{n\to\infty}\frac1n
 \log\sup_{u\in\pi(\Lambda_Q)}s_n(F,\varepsilon,\pi^{-1}(u)).
\end{align}
This is the precise input-relative topological quantity used below. We do not use the informal subtraction \(h_{\mathrm{top}}(F)-h_{\mathrm{top}}(\sigma)\), particularly when the latter is infinite.

The state-trajectory space is
\begin{align*}
 \mathcal T_Q=\{(\varphi(k,x,u))_{k\geq0}:(u,x)\in\Lambda_Q\}
       \subset Q^{\mathbb N_0}.
\end{align*}
With its product topology and left shift \(S\), it is a compact dynamical system: the trajectory map from \(\Lambda_Q\) is continuous and intertwines \(F\) and \(S\). A compatible metric on \(\mathcal T_Q\) is
\(\sum_{j\geq0}2^{-j-1}\min\{1,d(x_j,y_j)\}\).
For an \(F\)-invariant probability measure \(\mu\), the conditional metric entropy used here is the usual
\(h_\mu(F\mid\pi^{-1}\mathcal B_{\mathcal U})\). This notation refers to the input factor of the measure-preserving lift, not to an optimized control entropy.

\subsection{Linear coverings and radial estimates}

\begin{lemma}[Linear Bowen coverings]\label{lem:matrix}
Let \(L\in GL(d,\mathbb R)\) and let \(B\subset\mathbb R^d\) be compact. For each \(\delta>0\), there are finite sets \(E_n\subset B\) such that every \(y\in B\) has \(c\in E_n\) satisfying
\begin{align*}
 \max_{0\leq k\leq n}\left\lVert L^k(y-c)\right\rVert<\delta,
 \qquad
 |E_n|\leq C_\delta(n+1)^{d(d-1)}
                 \exp(n\chi^+(L)).
\end{align*}
The constant may depend on \(L,B,\delta\), but not on \(n\).
\end{lemma}
\begin{proof}
Work in real Jordan coordinates and group the real generalized eigenspaces according to eigenvalue modulus \(r>0\); a nonreal conjugate pair has its real dimension. On a group of dimension \(d_r\), the Jordan formula and equivalence of norms give
\begin{align*}
 \max_{0\leq k\leq n}\left\lVert L^kw\right\rVert
 \leq C(n+1)^{d-1}\max\{1,r^n\}\lVert w\rVert.
\end{align*}
Cover a fixed coordinate box containing \(B\) by half-open rectangular cells. In each coordinate of the group of modulus \(r\), take side length
\(\delta/[C'(n+1)^{d-1}\max\{1,r^n\}]\), where \(C'\) is large enough that two points in a cell have Bowen distance less than \(\delta\). Choose one point of \(B\) in each cell meeting \(B\). The number of cells is at most
\begin{align*}
 C_\delta(n+1)^{d(d-1)}
 \prod_r\max\{1,r^n\}^{d_r}.
\end{align*}
The product is \(\exp(n\chi^+(L))\). This proves the assertion, including nontrivial Jordan blocks and unit-modulus eigenvalues.
\end{proof}

\begin{lemma}[Radial decay]\label{lem:radial}
Suppose \(R:\mathbb R\to\mathbb R\) is an increasing \(C^\infty\) diffeomorphism, \(R(0)=0\), and \(0<R(s)<s\) for \(0<s\leq1\). Write \(\rho=R'(0)\), so \(0<\rho\leq1\). Then \(R^n(1)\downarrow0\), uniformly bounding \(R^n(s)\) for \(s\in[0,1]\). For every \(0<\eta<\rho\) there is \(C_\eta<\infty\) such that
\begin{align}\label{eq:radial-upper}
 0<\frac{R^n(s)}s\leq C_\eta(\rho+\eta)^n
       \quad(0<s\leq1,\ n\geq0).
\end{align}
For every \(s_0\in(0,1]\) there is \(c_{\eta,s_0}>0\) such that
\begin{align}\label{eq:radial-lower}
 \frac{R^n(s)}s\geq c_{\eta,s_0}(\rho-\eta)^n
       \quad(s_0\leq s\leq1,\ n\geq0).
\end{align}
The family \(\{R^n|_{[0,1]}:n\geq0\}\) is equicontinuous.
\end{lemma}
\begin{proof}
The limit of the decreasing sequence \(R^n(1)\) is a fixed point of \(R\) in \([0,1]\), hence is zero. Monotonicity gives \(0\leq R^n(s)\leq R^n(1)\). Choose \(a>0\) such that
\(\rho-\eta\leq R(t)/t\leq\rho+\eta\) for \(0<t\leq a\), and then \(N\) with \(R^N(1)\leq a\). In the product
\begin{align*}
 \frac{R^n(s)}s
 =\prod_{j=0}^{n-1}\frac{R(R^j(s))}{R^j(s)},
\end{align*}
all factors with \(j\geq N\) have the indicated bounds. The finitely many initial products have a uniform upper bound for \(0<s\leq1\), using their continuous extensions at zero. They have a positive minimum on \([s_0,1]\). Absorbing these finite products proves \eqref{eq:radial-upper} and \eqref{eq:radial-lower}, also for \(n<N\).

For equicontinuity, given \(\varepsilon>0\), choose \(N\) with \(R^N(1)<\varepsilon\). All maps with index at least \(N\) have range in \([0,\varepsilon)\). Uniform continuity of the finitely many preceding maps supplies a single initial-distance bound for the remaining indices.
\end{proof}

\subsection{Finite-horizon catalogues under uniform convergence}

The following estimate will allow the tracking upper bound to survive nonlinear perturbations. It uses exact controls for only a fraction of the requested tracking horizon.

\begin{proposition}[Shortening the exact-control horizon]\label{prop:shortening}
Let \(X\) be a normed vector space, let \(U\) be compact, let \(f:X\times U\to X\) be continuous, and let \(Q\subset X\) be compact and controlled invariant. Suppose that \(p\in Q\) and
\begin{align}\label{eq:uniform-collapse}
 b_n:=\sup\{\|\varphi(n,x,u)-p\|:(u,x)\in\Lambda_Q\}
 \longrightarrow0.
\end{align}
If \(H=h_{\mathrm{inv}}(Q,Q)<\infty\), then
\begin{align}\label{eq:tracking-exact-bound}
 h_{\mathrm{out}}(Q,Q)\leq h_{\mathrm{inv}}^+(Q,Q)\leq H.
\end{align}
Suppose in addition that \(b_n\leq Ca^n\) for constants \(C>0\), \(0<a<1\), and that all maps \(f_v\) are \(L\)-Lipschitz on the same ball about \(p\), with \(L>1\). Then
\begin{align}\label{eq:shortening-rate}
 h_{\mathrm{inv}}^+(Q,Q)
 \leq \frac{\log L}{\log(L/a)}\,H.
\end{align}
The same upper bounds hold after restricting the initial set to any compact subset of \(Q\).
\end{proposition}

\begin{proof}
Fix a sufficiently small \(\varepsilon>0\). Choose \(N\) so that \(2b_k<\varepsilon/2\) for every \(k\geq N\). Uniform continuity of the finitely many solution maps on \(Q\times U^N\) gives \(\delta>0\) such that two initial states at distance less than \(\delta\), subjected to the same input prefix, have distance less than \(\varepsilon/2\) at times \(0,\ldots,N\). Cover \(Q\) by \(C_\varepsilon\) sets of diameter less than \(\delta\).

Take an exact spanning catalogue of horizon \(m\geq N\). For each catalogue word and each covering set that contains an initial state feasible for this word through time \(m\), choose one such state as a reference. Extend its word to an infinite feasible input by controlled invariance. The extension may depend on the chosen reference; the number of reference-input pairs is at most
\begin{align}\label{eq:refinement-count}
 C_\varepsilon r_{\mathrm{inv}}(m,Q,Q).
\end{align}
Every initial state can be assigned a pair sharing its original catalogue word and its covering set. The actual and reference trajectories have distance less than \(\varepsilon/2\) through time \(N\). From time \(N\) through time \(m\), both trajectories are feasible, and their distance is at most \(2b_k<\varepsilon/2\). Here the bound for infinite feasible trajectories applies to each finite feasible prefix because that prefix can be continued in \(Q\). Taking \(m=n\), then normalized logarithms and the prescribed limits, proves \eqref{eq:tracking-exact-bound}.

For the sharper estimate let the common Lipschitz ball have radius \(r\). Decrease \(\varepsilon\) so that \(\varepsilon<r/4\), and increase \(N\) so that \(b_k<r/4\) for \(k\geq N\). Put
\begin{align}\label{eq:theta}
 \theta=\frac{\log L}{\log L-\log a}\in(0,1),
 \qquad m=\lceil\theta n+c_\varepsilon\rceil,
\end{align}
where \(c_\varepsilon>0\) is chosen so that
\(2C\exp[-c_\varepsilon\log(L/a)]<\varepsilon/2\).
For large \(n\), \(N\leq m\leq n\). Use the refined horizon-\(m\) catalogue just constructed. At time \(m\), the two assigned trajectories have distance at most \(2Ca^m\). For \(m\leq k\leq n\), iteration of the Lipschitz bound gives
\begin{align}\label{eq:tail-tracking}
 \|\varphi(k,x,u)-\varphi(k,a_i,u)\|
 \leq 2Ca^m L^{k-m}
 \leq 2Ca^m L^{n-m}<\varepsilon/2.
\end{align}
This iteration is justified inductively: the reference remains within \(r/4\) of \(p\), and the asserted error puts the actual state inside the ball of radius \(r/2\); hence both states lie in the common Lipschitz domain at the next step. The last inequality follows from \eqref{eq:theta}, since
\(n\log L-m\log(L/a)\leq-c_\varepsilon\log(L/a)\).
Consequently
\begin{align}
 r_{\mathrm{inv}}^+(n,\varepsilon,Q,Q)
 \leq C_\varepsilon
 r_{\mathrm{inv}}(\lceil\theta n+c_\varepsilon\rceil,Q,Q).
\end{align}
Because \(H<\infty\), the last spanning number is finite for all sufficiently large horizons. Dividing its logarithm by \(n\), using \(m/n\to\theta\), and taking the upper limit gives the bound \(\theta H\). Finally let \(\varepsilon\downarrow0\). Restriction of the initial set only decreases all the counts.
\end{proof}

\section{Entropy formulas on a solid cone}\label{sec:cone}

For this section fix the following assumptions:
\begin{enumerate}
\item \(A\in GL(d,\mathbb R)\), \(B\subset\mathbb R^d\) is compact and convex, and \(0\in\operatorname{int} B\);
\item \(U\subset\mathbb R^d\) is nonempty and compact, and
\begin{align}\label{eq:viability}
 \forall y\in B\quad\exists v\in U:\ Ay+v\in B;
\end{align}
\item \(R\) satisfies the assumptions of \cref{lem:radial}, and \(\rho=R'(0)\).
\end{enumerate}
Write
\begin{align*}
 t(s)=
 \begin{cases}R(s)/s,&s\ne0,\\ \rho,&s=0,\end{cases}
 \quad
 Q_B=\{(s,sy):0\leq s\leq1,\ y\in B\},
 \quad p=(0,0).
\end{align*}
We use Euclidean distances in the ambient space \(\mathbb R^{d+1}\). The controlled maps are
\begin{align}\label{eq:cone-system}
 f((s,z),v)=(R(s),t(s)(Az+sv)),\qquad v\in U.
\end{align}
The normalized affine system is \(g(y,v)=Ay+v\); superscripts \(f,g\) distinguish their spanning numbers when needed.

\Needspace{6\baselineskip}
The statement of \cref{thm:cone} is proved in the next four subsections.


The positive-volume hypothesis in part (ii) is a lower-bound hypothesis, not a controllability condition. The upper estimates hold even on the whole cone. Condition \eqref{eq:margin} is used only for the exact upper bound. It is sufficient and is not asserted to be necessary; the finite digit systems below attain that bound without satisfying this condition.

\subsection{Smoothness, normalization, and exact transference}

\begin{proof}[Proof of \cref{thm:cone}(i)]
The identity \(t(s)=\int_0^1R'(\theta s)\,d\theta\) shows that \(t\) is smooth and strictly positive on \(\mathbb R\). For fixed \(v\), the inverse of \(f_v\) at \((s',z')\) is
\begin{align*}
 s=R^{-1}(s'),\qquad z=A^{-1}\bigl(z'/t(s)-sv\bigr).
\end{align*}
This is globally defined and smooth. The set \(Q_B\) is the convex hull of \(p\) and \(\{1\}\times B\), so is compact and convex; since \(B\) has interior, this hull has nonempty interior in \(\mathbb R^{d+1}\).

For an initial state \((s,sy)\) with \(s>0\), let \(s_k=R^k(s)\). Direct substitution gives
\begin{align}\label{eq:normalized}
 \varphi(k,(s,sy),u)=(s_k,s_k y_k),\quad
 y_k=A^k y+b_k(u),\quad
 b_k(u)=\sum_{j=0}^{k-1}A^{k-1-j}u_j.
\end{align}
Here \(b_0=0\). Thus this state trajectory belongs to \(Q_B\) exactly when \(y_k\in B\). Recursive use of \eqref{eq:viability} gives an infinite viable normalized control from every \(y\in B\). The apex is fixed under every input.

Any finite catalogue spanning \(B\) for \(g\) therefore spans every positive radial slice and the apex for \(f\). Conversely, any catalogue spanning \(Q_B\) must span the slice \(\{(1,y):y\in B\}\); normalization shows that precisely the same catalogue spans \(B\) for \(g\). Existence of a finite catalogue in one problem is equivalent to its existence in the other, and minimizing proves \eqref{eq:transference}.
\end{proof}

\subsection{Finite outer and tracking catalogues}

\begin{proof}[Upper bounds in \cref{thm:cone}(ii)]
Fix \(\varepsilon>0\) and \(0<\eta<\rho\), and put \(\alpha=\rho+\eta\). Apply \cref{lem:matrix} to \(L=\alpha A\) and a radius \(\delta>0\) to be specified. For each of the finitely many centers \(c\in E_n\subset B\), choose once and for all an infinite normalized control \(u^c\) feasible from \(c\). Denote its normalized trajectory by \(c_k\in B\).

For \(y\in B\), choose \(c\in E_n\) with Bowen distance less than \(\delta\). The trajectories from \((s,sy)\) and \((s,sc)\) under the same \(u^c\) differ by
\begin{align*}
 (0,R^k(s)A^k(y-c)).
\end{align*}
By \eqref{eq:radial-upper}, \(s\leq1\), and the covering property, its norm is less than \(C_\eta\delta\) for \(0\leq k\leq n\). Choose \(\delta<\varepsilon/(2C_\eta)\). The reference trajectory from \((s,sc)\) is in \(Q_B\), so the actual trajectory is in \(Q_B^\varepsilon\). The same single control \(u^c\) works for all radial values associated with that center; the spanning catalogue is finite. Consequently
\begin{align*}
 \limsup_{n\to\infty}\frac1n
 \log r_{\mathrm{inv}}^{f}(n,Q_B,Q_B^\varepsilon)
 \leq\chi^+((\rho+\eta)A).
\end{align*}

For tracking, radial reference states must also belong to a \emph{finite} catalogue. Put \(C_B=\max_{y\in B}\left\lVert (1,y)\right\rVert\). Equicontinuity in \cref{lem:radial} and compactness of \([0,1]\) give a finite set \(T\subset[0,1]\), independent of \(n\), such that every \(s\in[0,1]\) has \(t\in T\) satisfying
\begin{align*}
 \sup_{k\geq0}|R^k(s)-R^k(t)|<\varepsilon/(2C_B).
\end{align*}
Use the finite pairs \(((t,tc),u^c)\), \(c\in E_n,\ t\in T\).
With \(t_k=R^k(t)\), their trajectory difference from \((s,sy)\) is
\begin{align*}
 (s_k-t_k,(s_k-t_k)c_k+s_kA^k(y-c)).
\end{align*}
Its norm is bounded by
\(C_B|s_k-t_k|+C_\eta\delta<\varepsilon\).
Every reference trajectory is feasible, including those with \(t=0\). Multiplication by the fixed number \(|T|\) does not change the exponential rate. Hence the tracking upper rate is also at most \(\chi^+((\rho+\eta)A)\), for every fixed \(\varepsilon>0\). Letting \(\eta\downarrow0\) gives the desired upper bound \(\chi^+(\rho A)\); the finite sum defining \(\chi^+\) is continuous under this scalar variation. Restricting the initial set to \(K\) preserves both bounds.
\end{proof}

\subsection{Unstable-volume lower bounds}

\begin{proof}[Lower bounds in \cref{thm:cone}(ii)]
There exists \(s_0>0\) such that
\begin{align*}
 K_0=K\cap\{(s,z):s\geq s_0\}
\end{align*}
has positive \((d+1)\)-dimensional volume: the increasing union of these truncated sets as \(s_0\downarrow0\) differs from \(K\) by a subset of \(\{s=0\}\), which has zero ambient volume.

Let \(E_+\) be the real sum of the generalized eigenspaces of \(A\) for which \(\rho|\lambda|>1\). Let \(E_-\) be the invariant complementary sum, let \(r=\dim E_+\), and write \(z=z_++z_-\). If \(r=0\), the claimed outer lower bound is simply nonnegativity. Suppose \(r>0\), and write \(A_+=A|_{E_+}\).

Fix a control \(u\). For arbitrary initial \(z\), \eqref{eq:normalized} gives the ambient last coordinate
\begin{align}\label{eq:ambient-final}
 z_n=a_n(s)A^n z+R^n(s)b_n(u),\qquad
 a_n(s)=R^n(s)/s.
\end{align}
Let \(W_n(u)\subset K_0\) consist of initial points whose trajectory stays in \(Q_B^\varepsilon\) at times \(0,\ldots,n\). Since \(Q_B^\varepsilon\) is bounded, its final \(E_+\)-coordinate belongs to a fixed ball, with radius depending on \(\varepsilon\) but not on \(n,u\). For each fixed \(s,z_-\), the map from \(z_+\) to that coordinate is an affine map with linear part \(a_n(s)A_+^n\). The \(r\)-dimensional volume of the corresponding section of \(W_n(u)\) is therefore at most
\begin{align*}
 C_\varepsilon\,a_n(s)^{-r}|\det A_+|^{-n}
 \leq C_{\varepsilon,\eta,s_0}
       \bigl((\rho-\eta)^r|\det A_+|\bigr)^{-n},
\end{align*}
where \eqref{eq:radial-lower} was used. This bound is unaffected by the translation, which may depend on \(u,s\).

The \(s,z_-\) projection of \(K_0\) is contained in a fixed bounded box. The change to the chosen linear coordinates has a constant nonzero Jacobian. Fubini's theorem consequently gives
\begin{align*}
 \operatorname{vol}_{d+1}(W_n(u))\leq
 C'_{\varepsilon,\eta,K}\bigl((\rho-\eta)^r|\det A_+|\bigr)^{-n}.
\end{align*}
The sets are measurable, since finite-time trajectories are continuous. If \(M\) controls span \(K\), their sets \(W_n(u)\) cover \(K_0\). Thus
\begin{align*}
 M\geq \frac{\operatorname{vol}_{d+1}(K_0)}{C'_{\varepsilon,\eta,K}}
          \bigl((\rho-\eta)^r|\det A_+|\bigr)^n.
\end{align*}
Taking the lower limit of normalized logarithms and then letting \(\eta\downarrow0\) yields
\begin{align*}
 \liminf_{n\to\infty}\frac1n
 \log r_{\mathrm{inv}}^{f}(n,K,Q_B^\varepsilon)
 \geq r\log\rho+\log|\det A_+|
 =\chi^+(\rho A).
\end{align*}
Together with the upper bound this proves an actual time limit for every \(\varepsilon>0\). Inequality \eqref{eq:outer-tracking} and the tracking upper bound give the same conclusion for tracking counts. The subsequent limit \(\varepsilon\downarrow0\) has the stated value.

For exact containment use instead the unstable subspace of \(A\), with \(|\lambda|>1\). From \(z_n/R^n(s)\in B\) we obtain
\begin{align*}
 A^nz+s b_n(u)\in sB\subset B.
\end{align*}
For fixed \(s,z_-\), the relevant linear part on the unstable subspace is now \(A_+^n\), without a radial factor. The same sectional-volume and Fubini argument gives
\begin{align*}
 \liminf_{n\to\infty}\frac1n
 \log r_{\mathrm{inv}}^{f}(n,K,Q_B)\geq\chi^+(A).
\end{align*}
This statement also covers infinite exact spanning numbers. It finishes part (ii).
\end{proof}

\subsection{The exact upper bound with an interior margin}

\begin{proof}[Proof of \cref{thm:cone}(iii)]
Choose \(\delta<\eta_0\) and a Bowen covering \(E_n\subset B\) for the matrix \(A\) from \cref{lem:matrix}. At each center \(c\), recursively use \eqref{eq:margin} to choose an infinite normalized control with
\begin{align*}
 \operatorname{dist}(c_k,\mathbb R^d\setminus B)\geq\eta_0
       \quad\text{for every }k\geq1.
\end{align*}
For any \(y\) assigned to \(c\), its normalized trajectory under that same control satisfies
\(\left\lVert y_k-c_k\right\rVert=\left\lVert A^k(y-c)\right\rVert<\delta\) for \(0\leq k\leq n\).
It follows that \(y_k\in B\) for \(k\geq1\), and \(y_0\in B\) already holds. Thus these finitely many controls span \(B\), and by exact transference they span \(Q_B\). Their number has upper exponential rate at most \(\chi^+(A)\). The exact lower bound just proved gives \eqref{eq:exact-spectral}.
\end{proof}

\begin{remark}[Where the assumptions enter]\label{rem:assumptions}
Convexity and \(0\in\operatorname{int} B\) make the constraint a solid convex cone; the proof uses boundedness for outer estimates and positive volume for lower estimates. No diagonalizability or hyperbolicity of \(A\) is needed. The smooth radial hypotheses ensure global invertibility, the estimates in \cref{lem:radial}, and uniform decay. Condition \eqref{eq:margin} is automatic if \(U\supset -AB\), since each normalized state can be reset to zero. It also holds for appropriate finite redundant alphabets. The theorem does not assert that an arbitrary viable compact affine system has finite exact invariance entropy without this condition.
\end{remark}

\begin{remark}[Why a Lie-group formula is not a direct substitute]\label{rem:lie}
In the author version of \cite[Definition 1]{CCS2021}, a discrete-time linear system on a Lie group has the form \(f_v(g)=a_v f_0(g)\), with \(f_0\) an automorphism. If all such maps fix one point \(p\), then \(a_v=p f_0(p)^{-1}\) is independent of \(v\), and all controlled maps coincide. The cone maps have a common fixed apex and are distinct for distinct \(v\). Thus they do not have this representation on any group structure that represents the same maps on the whole state space. The normalization \(y=z/s\) is defined only for \(s>0\); deleting the apex removes compactness of the viable constraint and does not justify transferring a compact-constraint entropy formula. This observation excludes that direct application, not every possible use of homogeneous-space theory.
\end{remark}

\subsection{Nonlinear perturbations with fixed first jet}

We prove \cref{thm:nonlinear}. Exact transference no longer holds for these systems: the normalized dynamics depends on the current radius. Its entropy can nevertheless be computed by combining finite blocks at small radii with one-dimensional sectional-volume estimates.

\begin{proof}[Proof of \cref{thm:nonlinear}]
Choose \(\delta\) with
\(\|P\|_{C_b^2}<\delta<\min\{\eta,q/2\}\).
Taylor's formula at the origin gives, for \(0<s\leq1\) and \(|y|\leq1\),
\begin{align}\label{eq:perturbation-small}
 |P_v(s,sy)|\leq\tfrac12\delta s^2(1+y^2)\leq\delta s^2,
 \qquad |\partial_zP_v(s,z)|\leq\delta\sqrt{s^2+z^2}.
\end{align}
The second inequality holds for all \((s,z)\); the global bound \(|\partial_zP_v|<\delta\) also holds.

\emph{Smoothness and viability.}
The first coordinate determines \(s\) uniquely from the image. For this \(s\), the map
\(z\mapsto qz+sv+P_v(s,z)\) has derivative at least \(q-\delta>0\), tends to the two respective infinities at the ends of the real line, and is therefore a bijection. The implicit function theorem gives a smooth inverse depending on \(s\) and the image coordinate. Positivity and smoothness of \(t\), already established above, prove global smooth invertibility.

At positive radius the normalized coordinate satisfies
\begin{align}\label{eq:nonlinear-normalized}
 y_{k+1}=qy_k+u_k+
 \frac{P_{u_k}(s_k,s_ky_k)}{s_k},\qquad
 s_k=R^k(s).
\end{align}
At any \(|y_k|\leq1\), choose \(u_k\) by \eqref{eq:nonlinear-margin}. The perturbation term has absolute value at most \(\delta s_k\leq\delta<\eta\), so \(|y_{k+1}|\leq1-\eta\). The apex is fixed under every input. This proves controlled invariance, with a uniform normalized interior margin. Every finite feasible trajectory satisfies \(0\leq s_k\leq R^k(1)\) and \(|z_k|\leq s_k\). Hence all its time-\(k\) states lie in \(R^k(1)Q\).

\emph{Exact upper bound.}
Fix an integer block length \(N\geq1\). Cover \([-1,1]\) by intervals of radius
\(\gamma_N=(\eta/4)q^{-N}\), with centers in \([-1,1]\), using at most \(C_\eta q^N\) intervals, where \(C_\eta\) is independent of \(N\). In the unperturbed affine system \(y^+=qy+v\), choose from each center a length-\(N\) word whose reference states at positive times all satisfy \(|c_k|\leq1-2\eta\). This is possible by recursive use of \eqref{eq:nonlinear-margin}.

Choose \(a_N\in(0,1]\) so small that
\begin{align}\label{eq:small-radius-block}
 \delta a_N\sum_{j=0}^{N-1}q^j\leq\eta/4.
\end{align}
For initial radius \(0<s\leq a_N\) and an initial ratio \(y\) assigned to a center, apply its unperturbed word to the perturbed system. As long as the perturbed ratios remain in \([-1,1]\), their errors relative to the unperturbed reference obey
\begin{align}
 |y_k-c_k|
 \leq q^k\gamma_N+\delta a_N\sum_{j=0}^{k-1}q^j
 \leq\eta/2\qquad(0\leq k\leq N).
\end{align}
This estimate proves its own feasibility premise by induction: the reference at each positive time has margin \(2\eta\), while the error is at most \(\eta/2\). Thus these same \(C_\eta q^N\) words work simultaneously for all radii at most \(a_N\). No radius-dependent infinite catalogue has been introduced.

Choose \(N_0\) with \(R^{N_0}(1)\leq a_N\). All words in \(U^{N_0}\) form a finite exact catalogue for this transient, by controlled invariance. Thereafter concatenate arbitrary words from the small-radius block catalogue. Each current radius is still at most \(a_N\), so the argument can be applied successively to every block. For \(n\geq N_0\), using a prefix of the final block if necessary, this gives
\begin{align}
 r_{\mathrm{inv}}(n,Q,Q)
 \leq |U|^{N_0}(C_\eta q^N)^{\lceil(n-N_0)/N\rceil}.
\end{align}
Consequently the upper entropy is at most
\(\log q+N^{-1}\log C_\eta\). Letting \(N\to\infty\) proves the exact upper bound \(\log q\).

\emph{Exact lower bound.}
For fixed \(s>0\) and a fixed input word \(u\), let \(I_n(s,u)\) be the set of \(z\in[-s,s]\) whose states satisfy \(Q\) through time \(n\). This set is an interval, possibly empty or a singleton: every vertical solution map is strictly increasing, and \(I_n(s,u)\) is the intersection of inverse images of vertical intervals. In particular every point between two members has a feasible trajectory through time \(n\).

Write \(Z_n(s,z,u)\) for the second state coordinate. Along this interval the chain rule yields
\begin{align}\label{eq:vertical-derivative}
 \partial_z Z_n(s,z,u)
 =\frac{R^n(s)}s
   \prod_{j=0}^{n-1}
       \bigl[q+\partial_zP_{u_j}(s_j,z_j)\bigr].
\end{align}
The product of the factors \(t(s_j)\) telescopes to \(R^n(s)/s\).
For any \(\zeta\in(0,q-1)\), uniform feasible convergence and \eqref{eq:perturbation-small} imply that every factor after a fixed time \(N_\zeta\) is at least \(q-\zeta\), independently of \(s,z,u\). The finitely many preceding factors are at least \(q-\delta\). There is therefore a constant \(c_\zeta>0\), independent of the word and of \(n,s\), such that
\begin{align}
 \partial_z Z_n(s,z,u)
 \geq c_\zeta\frac{R^n(s)}s(q-\zeta)^n
 \quad\text{on }I_n(s,u).
\end{align}
The image interval at time \(n\) has length at most \(2R^n(s)\). Integrating the derivative between any two points of \(I_n(s,u)\) proves
\begin{align}\label{eq:exact-section-nonlinear}
 |I_n(s,u)|\leq \frac{2s}{c_\zeta}(q-\zeta)^{-n}.
\end{align}
The feasible set of one word is closed and hence measurable. Fubini's theorem, applied over \(0<s\leq1\), bounds its area in \(Q\) by a constant times \((q-\zeta)^{-n}\). The slice \(s=0\) has zero area. Covering a compact \(K\subset Q\) of positive area consequently requires at least \(c_K(q-\zeta)^n\) words. Taking the lower time limit and then \(\zeta\downarrow0\), and combining with the upper bound, proves the first line of \eqref{eq:nonlinear-entropies}, including existence of its time limit.

\emph{Outer lower bound.}
Choose \(s_0>0\) so that \(K_0=K\cap\{s\geq s_0\}\) has positive area. Fix small \(\varepsilon>0\). Replace \(I_n(s,u)\) by the initial interval of points in \([-s,s]\) whose trajectory remains in \(Q^\varepsilon\) through time \(n\). This is still an interval: \(Q^\varepsilon\) is convex, all its vertical sections are intervals, and all vertical solution maps are increasing. For an outer-feasible state with radial coordinate \(s_j\), comparison with a point \((r,w)\in Q\) at distance less than \(\varepsilon\) gives
\begin{align}
 |z_j|<s_j+2\varepsilon.
\end{align}
After a time depending only on \(\varepsilon\), \(s_j\leq R^j(1)\leq\varepsilon\), and hence \(|z_j|\leq3\varepsilon\). By \eqref{eq:perturbation-small}, the late factors in \eqref{eq:vertical-derivative} are bounded below by
\(q-\omega_\varepsilon\), where
\(\omega_\varepsilon=\delta\sqrt{10}\,\varepsilon\to0\).
Take \(\varepsilon\) small enough that this number is positive. Absorbing the finitely many earlier factors into \(c_\varepsilon>0\), the derivative on the entire outer-feasible initial interval is at least
\begin{align}
 c_\varepsilon\frac{R^n(s)}s(q-\omega_\varepsilon)^n.
\end{align}
Its time-\(n\) image has length at most \(2(1+2\varepsilon)\). For \(s\geq s_0\), the radial lower bound in \cref{lem:radial} then bounds its initial length by
\begin{align}
 C_{\varepsilon,\xi,s_0}
 \bigl[(\rho-\xi)(q-\omega_\varepsilon)\bigr]^{-n},
 \qquad 0<\xi<\rho.
\end{align}
The mean-value estimate remains valid for open or half-open intervals by applying it between any two of their points. Fubini on \(K_0\), followed by the covering argument, gives a lower time rate at least
\(\log[(\rho-\xi)(q-\omega_\varepsilon)]\).
Nonnegativity supplies the maximum with zero. Let \(\xi\downarrow0\) and then \(\varepsilon\downarrow0\), in that order after the time limit, to obtain
\begin{align}\label{eq:nonlinear-outer-lower}
 h_{\mathrm{out}}(K,Q)\geq\max\{0,\log(\rho q)\}.
\end{align}

\emph{Tracking upper bound.}
If \(\rho=1\), uniform feasible convergence and \cref{prop:shortening} give
\(h_{\mathrm{inv}}^+(Q,Q)\leq h_{\mathrm{inv}}(Q,Q)=\log q\), which matches \eqref{eq:nonlinear-outer-lower}.
Suppose \(0<\rho<1\). For any \(a\in(\rho,1)\), \cref{lem:radial} and feasibility give a uniform bound \(b_n\leq C_a a^n\).
At the apex,
\begin{align}\label{eq:apex-derivative}
 Df_v(p)=
 \begin{pmatrix}
 \rho&0\\
 \rho v&\rho q
 \end{pmatrix}.
\end{align}
For \(M>0\), equip \(\mathbb R^2\) with the equivalent norm
\(\|(s,z)\|_M=\max\{M|s|,|z|\}\).
The corresponding operator norm of \eqref{eq:apex-derivative} is at most
\(\rho q+\rho\max_{v\in U}|v|/M\).
For any \(L>\max\{1,\rho q\}\), choose \(M\) large and then a common convex neighborhood of \(p\) small enough that all \(f_v\) are \(L\)-Lipschitz there. Finiteness of \(U\) makes the neighborhood uniform. Equivalent norms leave these entropy definitions unchanged, since each metric tolerance can be bounded by a constant multiple of the other.

Applying \cref{prop:shortening} in this norm yields
\begin{align}
 h_{\mathrm{inv}}^+(Q,Q)
 \leq \log q\,\frac{\log L}{\log L-\log a}.
\end{align}
Let \(a\downarrow\rho\) and \(L\downarrow\max\{1,\rho q\}\).
If \(\rho q>1\), the right-hand side tends to \(\log(\rho q)\); if \(\rho q\leq1\), it tends to zero. This includes the critical case \(\rho q=1\). Restricting the initial set to \(K\) and using \(h_{\mathrm{out}}\leq h_{\mathrm{inv}}^+\) completes the proof.
\end{proof}

\begin{remark}[What the nonlinear calculation does and does not preserve]\label{rem:nonlinear-scope}
The exact entropy is determined by the limiting normalized expansion \(q\). Its upper bound uses finite blocks that work uniformly for a whole radial interval, whereas its lower bound uses derivatives on vertical intervals. The outer lower bound additionally uses the smallness of the derivative perturbation near the apex; the tracking upper bound uses \cref{prop:shortening}. Exact finite-horizon transference and the digit tiling are not claimed for the perturbed system. Nor do we assert a common fixed-\(\varepsilon\) time limit for its outer or tracking counts. The entropy identities use exactly the iterated limits in \cref{sec:definitions}.
\end{remark}

\section{Finite alphabets and equal positive entropies}\label{sec:examples}

\subsection{Exact counts in arbitrary dimension}

For integers \(q_1,\ldots,q_d\geq2\), set
\begin{align*}
 U_{q_i}=\{q_i-1-2j:0\leq j<q_i\},\quad
 A=\operatorname{diag}(q_1,\ldots,q_d),\quad
 U=\prod_{i=1}^d U_{q_i},\quad M=\prod_{i=1}^d q_i.
\end{align*}
Take \(B=[-1,1]^d\), \(Q=Q_B\), and
\begin{align}\label{eq:K}
 K_d=[1/2,3/4]\times[-1/4,1/4]^d.
\end{align}
This compact set is contained in \(\operatorname{int} Q\) and has positive ambient volume. In particular, \(|z_i|\leq1/4<1/2\leq s<1\) on \(K_d\).

\begin{theorem}[Finite digit family]\label{thm:digits}
For every radial map \(R\) satisfying \cref{lem:radial}, the preceding finite-alphabet cone system satisfies, for all integers \(n\geq0\),
\begin{align}\label{eq:digit-counts}
 r_{\mathrm{inv}}(n,Q,Q)=M^n,\qquad
 2^{-d}M^n\leq r_{\mathrm{inv}}(n,K_d,Q)\leq M^n.
\end{align}
Consequently
\begin{align*}
 h_{\mathrm{inv}}(K_d,Q)=h_{\mathrm{inv}}(Q,Q)=\sum_{i=1}^d\log q_i,
\end{align*}
and
\begin{align}\label{eq:digit-outer}
 h_{\mathrm{out}}(K_d,Q)=h_{\mathrm{out}}(Q,Q)
 =h_{\mathrm{inv}}^{+}(K_d,Q)=h_{\mathrm{inv}}^{+}(Q,Q)
 =\sum_{i=1}^d\max\{0,\log(\rho q_i)\}.
\end{align}
For any two initial states \(x,y\in Q\) and any two controls \(u,v\) respectively feasible for all nonnegative times,
\begin{align}\label{eq:uniform-digit}
 \left\lVert \varphi(n,x,u)-\varphi(n,y,v)\right\rVert
 \leq 2\sqrt d\,R^n(1)\longrightarrow0.
\end{align}
All individual controlled maps are global smooth diffeomorphisms.
\end{theorem}
\begin{proof}
For one coordinate with expansion \(q\), the intervals
\begin{align*}
 I_v=\left[\frac{-1-v}{q},\frac{1-v}{q}\right],
       \quad v\in U_q,
\end{align*}
tile \([-1,1]\) with disjoint interiors. The map \(y\mapsto qy+v\) sends \(I_v\) onto \([-1,1]\), and its inverse branch sends all of \([-1,1]\) into \(I_v\). Hence every point admits an infinite viable sequence, and every word of length \(n\) has a viable interval of length \(2q^{-n}\). More explicitly, composing the inverse branches maps \([-1,1]\) to the set viable for that word; all intermediate inverse images remain in \([-1,1]\).

Products give one viable box for each length-\(n\) word over \(U\), of volume \(2^dM^{-n}\). The \(M^n\) words span the entire cube \(B\). Conversely, for any fixed word the final-time condition alone confines its viable initial points to an affine preimage of \(B\) of volume \(2^dM^{-n}\). Covering \(B\), of volume \(2^d\), requires at least \(M^n\) words. Transference proves the first identity in \eqref{eq:digit-counts}.

To obtain the lower bound for \(K_d\), restrict any spanning family to the slice \(s=1/2\). Its normalized initial coordinates contain exactly \([-1/2,1/2]^d\), of \(d\)-volume \(1\). One control word covers at most \(2^dM^{-n}\) of that slice in normalized coordinates. At least \(2^{-d}M^n\) words are required. The upper bound follows by restricting a catalogue for \(Q\). This gives the exact entropy formula without an interior margin assumption. Formula \eqref{eq:digit-outer} follows from \cref{thm:cone}(ii).

Every feasible state at time \(n\) belongs to
\(\{(s,z):0\leq s\leq R^n(1),\ |z_i|\leq s\}=R^n(1)Q\).
The diameter of \(Q\) is \(2\sqrt d\): its vertices are the apex and the vertices \((1,\pm1,\ldots,\pm1)\), and the largest distance between them is \(2\sqrt d\). A convex hull has the same diameter as the supremum of the distances between its generating vertices. This proves \eqref{eq:uniform-digit}. Smooth invertibility was proved in \cref{thm:cone}(i).
\end{proof}

\begin{remark}[The planar linear-radial system]
With \(d=1\) and \(R(s)=\rho s\), \(0<\rho<1\), the system is
\begin{align*}
 f((s,z),v)=(\rho s,\rho(qz+sv)),\quad v\in U_q,
 \quad Q=\{0\leq s\leq1,\ |z|\leq s\}.
\end{align*}
Its fixed-input maps are linear isomorphisms with determinant \(\rho^2q\); its dependence on state and input is bilinear. The counts reduce to
\begin{align*}
 r_{\mathrm{inv}}(n,Q,Q)=q^n,\qquad
 q^n/2\leq r_{\mathrm{inv}}(n,K_1,Q)\leq q^n.
\end{align*}
The exact rate is \(\log q\), the outer and tracking rates are \(\max\{0,\log(q\rho)\}\), and the distance of any two feasible trajectories is at most \(2\rho^n\). In particular both control entropies are positive for \(1/q<\rho<1\).
\end{remark}

\subsection{Subexponential radial decay removes the entropy gap}

\begin{corollary}[Equal positive entropies]\label{cor:equal}
Let \(R(s)=\operatorname{arsinh}s\), \(d=1\), and \(q=2\), with \(U_2=\{-1,1\}\). On the same solid triangle \(Q\) and interior rectangle \(K_1\), every controlled map is a global smooth diffeomorphism and
\begin{align}\label{eq:equal}
 h_{\mathrm{inv}}(K_1,Q)=h_{\mathrm{out}}(K_1,Q)=h_{\mathrm{inv}}^{+}(K_1,Q)=\log2.
\end{align}
The same equalities hold for \(K_1\) replaced by \(Q\). All pairs of feasible state trajectories are asymptotic. Moreover,
\begin{align*}
 R^n(1)\sim\sqrt{3/n}.
\end{align*}
\end{corollary}
\begin{proof}
The function \(R=\operatorname{arsinh}\) is a smooth increasing bijection of \(\mathbb R\) with smooth inverse \(\sinh\), satisfies \(0<R(s)<s\) for \(s>0\), and has \(R'(0)=1\). Thus \cref{thm:digits} applies with \(\rho=1\), proving \eqref{eq:equal} and asymptoticity.

For the rate, write \(s_n=R^n(1)\). Since
\(R(s)=s-s^3/6+O(s^5)\) at zero, and \(s_n\downarrow0\),
\begin{align*}
 s_{n+1}^{-2}-s_n^{-2}
 =s_n^{-2}\bigl[(1-s_n^2/6+O(s_n^4))^{-2}-1\bigr]
 \longrightarrow1/3.
\end{align*}
Telescoping and taking Ces\`aro averages gives \(s_n^{-2}/n\to1/3\), which proves the asserted asymptotic formula.
\end{proof}

The equality case shows that a discrepancy between exact and neighborhood constraints is not responsible for the failure of a state Li--Yorke conclusion. A fixed positive ambient tolerance still requires exponentially many control signals. Nevertheless every signal that is infinitely feasible produces a state converging to the apex.

\subsection{A nonlinear equality family}\label{sec:joint-proof}

\begin{proof}[Proof of \cref{cor:joint}]
For \(q=2\) and \(U=\{-3/2,0,3/2\}\), the intervals of points for which one of \(|2y+v|\) is at most \(3/4\) cover \([-1,1]\). Equivalently, the nearest of the three centers \(-3/4,0,3/4\) is at distance at most \(3/8\) from every \(y\in[-1,1]\). Thus \eqref{eq:nonlinear-margin} holds with \(\eta=1/8\).

The function \(\tanh(s)\tanh(z)\) and its first two derivatives are bounded; its value and first derivative at the origin vanish. The absolute values of its two pure second derivatives are at most \(2\), and that of its mixed derivative is at most \(1\), so its Hessian operator norm is at most \(3\). In particular the bound \(4\max_v|e_v|<1/16<\eta\) dominates the required \(C_b^2\) norm. \Cref{thm:nonlinear} therefore gives the three entropy values and uniform feasible convergence. The diameter of \(Q\) is \(2\), giving the pair-distance bound. The radial asymptotic is proved in \cref{cor:equal}.

On \(Q\), the function \(V(s,z)=s\) vanishes only at the apex and satisfies
\(V(f_v(s,z))=R(V(s,z))\) for every feasible step. Apply \cref{thm:lyapunov} with \(\psi=R|_{[0,1]}\). Since \(|U|=3\), its lift, fiber, state-factor, and measure conclusions give all remaining assertions.
\end{proof}

Taking at least one \(e_v\ne0\) makes the vertical controlled map nonlinear. More generally the entire perturbation ball in \cref{thm:nonlinear}, with this fixed alphabet and \(R=\operatorname{arsinh}\), has the same entropy equalities and the same Lyapunov conclusions. Redundancy of the alphabet supplies the interior margin that allows this perturbation argument; the minimal digit alphabet is used for exact counting, not for this persistence statement.

\subsection{Prescribing the two control entropy rates}

\begin{corollary}[Arbitrary ordered rates]\label{cor:rates}
For every \(a>0\) and \(b\in[0,a]\), there is a planar system of the form \eqref{eq:cone-system} with a finite input alphabet, global smooth invertible controlled maps, the triangle \(Q\), and the rectangle \(K_1\), such that
\begin{align*}
 h_{\mathrm{inv}}(K_1,Q)=h_{\mathrm{inv}}(Q,Q)=a,\qquad
 h_{\mathrm{out}}(K_1,Q)=h_{\mathrm{out}}(Q,Q)=b.
\end{align*}
Its tracking entropies also equal \(b\), and all pairs of feasible state trajectories are asymptotic.
\end{corollary}
\begin{proof}
Put \(L=e^a>1\), \(m=\lceil L\rceil+2\), and
\begin{align*}
 c_j=-1+\frac{2j}{m-1},\qquad
 U=\{-Lc_j:0\leq j<m\},\qquad A=L,\quad B=[-1,1].
\end{align*}
For every \(y\in B\), a nearest center satisfies \(|y-c_j|\leq1/(m-1)\), so
\begin{align*}
 |Ly-Lc_j|\leq\frac L{m-1}<1.
\end{align*}
Thus \eqref{eq:margin} holds with \(\eta_0=1-L/(m-1)>0\). Exact entropy equals \(\log L=a\) by \cref{thm:cone}(iii). If \(b<a\), choose \(R(s)=e^{b-a}s\). Then \(\rho L=e^b\) and the outer and tracking formula is \(b\), including the threshold case \(b=0\). If \(b=a\), choose \(R(s)=\operatorname{arsinh}s\); then \(\rho=1\) gives the same conclusion. Uniform radial decay and feasibility place all trajectories in \(R^n(1)Q\), proving asymptoticity.
\end{proof}

This realizes the full range of ordered finite values in dimension two, with a finite alphabet depending on \(a\). It is not an assertion that the chosen alphabet size is minimal. Exact entropy cannot be smaller than outer entropy by definition, so the ordering in the corollary is the relevant restriction.

\section{A Lyapunov condition and entropy separation}\label{sec:separation}

\subsection{A system-level obstruction}

The next result uses a one-step inequality for a scalar function. It does not assume shrinking trajectory diameters or any positive invariance entropy. The cone formula supplies positive control entropy separately, under explicit algebraic and viability assumptions.

\begin{theorem}[Common Lyapunov obstruction]\label{thm:lyapunov}
Let \(X,U,f,Q\) be as in \cref{sec:definitions}, and suppose \(p\in Q\) satisfies \(f(p,v)=p\) for every \(v\in U\). Suppose there is a continuous function \(V:Q\to[0,L]\), with \(L=\max_QV\), such that \(V^{-1}(0)=\{p\}\), and a continuous nondecreasing function \(\psi:[0,L]\to[0,L]\) with \(\psi(0)=0\) and \(\psi(t)<t\) for \(t>0\), satisfying
\begin{align}\label{eq:lyapunov}
 V(f(x,v))\leq\psi(V(x))
 \quad\text{whenever }x\in Q,\ v\in U,\ f(x,v)\in Q.
\end{align}
Then:
\begin{enumerate}
\item all infinite feasible trajectories converge to \(p\), uniformly over their initial states and their controls; every two such trajectories are asymptotic;
\item the state-trajectory shift has \(h_{\mathrm{top}}(S)=0\), and the uniform fiber entropy in \eqref{eq:fiber-definition} satisfies \(h_{\mathrm{fib}}(F\mid\pi)=0\);
\item \(\pi(\Lambda_Q)=\mathcal U\) and
\begin{align*}
 h_{\mathrm{top}}(F)=h_{\mathrm{top}}(\sigma:\mathcal U\to\mathcal U)
 =\begin{cases}\log m,&|U|=m<\infty,\\+\infty,&U\text{ infinite};\end{cases}
\end{align*}
\item every \(F\)-invariant Borel probability measure is supported on \(\mathcal U\times\{p\}\). On that set \(\pi\) is a conjugacy to the input shift. In particular
\begin{align*}
 h_\mu(F\mid\pi^{-1}\mathcal B_{\mathcal U})=0;
\end{align*}
\item the eventual core is
\begin{align*}
 \bigcap_{n\geq0}F^n(\Lambda_Q)=\mathcal U\times\{p\}.
\end{align*}
If \(Q\ne\{p\}\), then \(F\) is not surjective.
\end{enumerate}
Every continuous self-map \(G:Q\to Q\) obtained from a feasible memoryless feedback has zero topological entropy. All feasible trajectories of any fixed causal controller, whether continuous or not and whether or not it uses memory, satisfy the asymptotic conclusion in (i).
\end{theorem}

\begin{proof}
If \(L=0\), then \(Q=\{p\}\), and all assertions follow directly; assume \(L>0\). The numbers \(\psi^n(L)\) decrease to a fixed point of \(\psi\). Continuity and strict decrease away from zero force this limit to be zero. Iterating \eqref{eq:lyapunov} along a feasible trajectory gives
\begin{align*}
 V(\varphi(n,x,u))\leq\psi^n(L).
\end{align*}
For any \(\varepsilon>0\), the compact set
\(\{x\in Q:d(x,p)\geq\varepsilon\}\), if nonempty, has a strictly positive minimum of \(V\). Thus there are numbers \(\beta_n\downarrow0\) with
\begin{align}\label{eq:beta}
 \sup_{(u,x)\in\Lambda_Q}d(\varphi(n,x,u),p)\leq\beta_n.
\end{align}
For example take the maximum distance to \(p\) over \(V^{-1}([0,\psi^n(L)])\).
This proves (i), including the quantification over different controls.

We give uniform covering arguments for (ii) and (iii) to avoid identifying different relative entropy conventions. Fix \(0<\varepsilon<1\) and choose \(N\) such that \(2\beta_k<\varepsilon/2\) for \(k\geq N\). The finitely many maps
\((x,u_0,\ldots,u_{N-1})\mapsto\varphi(k,x,u)\), \(0\leq k<N\), are uniformly continuous on the compact domain \(Q\times U^N\). This remains true even though intermediate states for some controls may leave \(Q\); each finite-time image of the compact domain is compact, and \(f\) is continuous on \(X\times U\). There exists \(\delta>0\), independent of the control prefix, such that
\begin{align}\label{eq:prefix-uc}
 d(x,y)<\delta\quad\Longrightarrow\quad
 \max_{0\leq k<N}d(\varphi(k,x,u),\varphi(k,y,u))<\varepsilon/2
\end{align}
for every common prefix.

Cover \(Q\) by finitely many sets of diameter less than \(\delta\), say \(C\) sets. Within one input fiber, two feasible initial states in the same such set have state distance less than \(\varepsilon/2\) up to time \(N-1\) by \eqref{eq:prefix-uc}. At later times their distance is less than \(\varepsilon/2\) by \eqref{eq:beta}. Their input coordinates are identical at all shifted times. Consequently every \((n,\varepsilon)\)-separated subset of every fiber has cardinality at most \(C\), uniformly in \(n\) and \(u\). This proves \(h_{\mathrm{fib}}(F\mid\pi)=0\).

For the state-trajectory system, \eqref{eq:beta} implies
\begin{align*}
 \sup_{\xi\in\mathcal T_Q}d_{\mathcal T_Q}(S^n\xi,(p,p,\ldots))
 \leq\beta_n\longrightarrow0.
\end{align*}
Any continuous map on a compact metric space whose iterates converge uniformly to a constant has zero entropy: for a given separation scale, all sufficiently late images have diameter smaller than that scale, while uniform continuity of the finitely many earlier iterates gives a finite cover of bounded Bowen diameter for every horizon. Applying this argument proves \(h_{\mathrm{top}}(S)=0\).

The set \(\mathcal U\times\{p\}\) belongs to \(\Lambda_Q\), and \(F\) restricted to it is conjugate to the full input shift. This proves surjectivity of \(\pi\) and the lower entropy bound in (iii). If \(|U|=m<\infty\), choose \(\ell\) so that \(2^{-\ell}<\varepsilon/2\). For large \(n\), group lift points by their first \(n+\ell\) input symbols and by a member of the above finite state cover. Points in one group have close initial state evolution until time \(N\), since their input prefixes agree and \(n+\ell\geq N\). Their later feasible state evolutions are close by \eqref{eq:beta}. At each shifted time \(0\leq k<n\), agreement of the input prefix gives input distance at most \(2^{-\ell}<\varepsilon/2\). Thus each group has Bowen diameter less than \(\varepsilon\), and
\begin{align*}
 s_n(F,\varepsilon)\leq C m^{n+\ell}.
\end{align*}
This proves the upper bound \(\log m\). If \(U\) is infinite, it contains finite subsets of every cardinality \(m\). Each corresponding full shift at the apex is a compact invariant subsystem of entropy \(\log m\). Letting \(m\to\infty\) proves that both full input entropy and total lift entropy are infinite.

For (iv), regard \(V\) also as a function on the lift. Invariance of \(\mu\) gives
\begin{align*}
 0=\int_{\Lambda_Q}(V-V\circ F)\,d\mu
 \geq\int_{\Lambda_Q}(V-\psi\circ V)\,d\mu\geq0.
\end{align*}
The last integrand is nonnegative and vanishes only when the state is \(p\). Hence \(\mu(\mathcal U\times\{p\})=1\). Projection and \(u\mapsto(u,p)\) are inverse measurable maps on this full-measure set and intertwine the transformations. The input factor therefore contains the whole state information modulo \(\mu\), and its conditional Kolmogorov--Sinai entropy is zero. Conversely, any invariant input measure lifts in this way.

For (v), every point of \(F^n(\Lambda_Q)\) has \(V\)-value at most \(\psi^n(L)\), so a point in the intersection has state \(p\). Conversely every \((u,p)\) admits arbitrarily long prehistories in the lift: prepend any \(n\) input values to \(u\) and keep the state at \(p\). This proves the core formula. Finally some \(x\in Q\) attains \(V(x)=L>0\); it has an infinite feasible control by controlled invariance. Its lift point cannot be in \(F(\Lambda_Q)\), whose \(V\)-values are at most \(\psi(L)<L\). Hence \(F\) is not onto.

A feasible controller generates only trajectories already covered by (i), regardless of how it makes its choices. If it induces a continuous self-map \(G\) of \(Q\), then \(G^n\to p\) uniformly; the same bounded-Bowen-cover argument used for \(S\) gives \(h_{\mathrm{top}}(G)=0\).
\end{proof}

For cone systems, take \(V(s,z)=s\), \(L=1\), and \(\psi=R|_{[0,1]}\). On \(Q_B\), zero radial coordinate occurs only at the apex and \eqref{eq:lyapunov} holds with equality. The preceding theorem therefore applies to every system in \cref{thm:cone}, independently of its spectral assumptions for positive entropy.

\subsection{The entire digit lift}

\begin{proposition}[Explicit coding of the true lift]\label{prop:coding}
For the digit system of \cref{thm:digits}, let \(\Sigma=U^{\mathbb N_0}\) and define
\begin{align*}
 Y_i(u)=-\sum_{j=0}^{\infty}\frac{u_{j,i}}{q_i^{j+1}},
       \qquad Y(u)=(Y_1(u),\ldots,Y_d(u)).
\end{align*}
Then
\begin{align}\label{eq:H}
 H:\Sigma\times[0,1]\longrightarrow\Lambda_Q,\qquad
 H(u,s)=(u,(s,sY(u)))
\end{align}
is a homeomorphism, and
\begin{align*}
 F\circ H=H\circ(\sigma\times R).
\end{align*}
In particular the total lift entropy is \(\log M\), its uniform input-fiber entropy is zero, its state-trajectory entropy is zero, and every invariant lift measure is \(H_*(\nu\times\delta_0)\) for a shift-invariant probability measure \(\nu\) on \(\Sigma\).
\end{proposition}
\begin{proof}
Since \(|u_{j,i}|\leq q_i-1\), the series converge uniformly and \(|Y_i(u)|\leq1\); \(Y\) is continuous in the product topology. The identity
\begin{align*}
 q_iY_i(u)+u_{0,i}=Y_i(\sigma u)
\end{align*}
shows that \(H(u,s)\) is feasible and proves equivariance.

Conversely, let \((u,(s,z))\in\Lambda_Q\) with \(s>0\), and put \(y=z/s\). The normalized trajectory remains in \([-1,1]^d\). For each coordinate, dividing the \(n\)-step equation by \(q_i^n\) gives
\begin{align*}
 y_i=-\sum_{j=0}^{n-1}\frac{u_{j,i}}{q_i^{j+1}}
           +q_i^{-n}y_{n,i}.
\end{align*}
The final term tends to zero. Thus \(y=Y(u)\). If \(s=0\), the state is the apex and is \(H(u,0)\). This proves surjectivity onto the \emph{entire} feasible lift, not merely onto a selected invariant subset. The input and radial coordinates recover \(u,s\) uniquely, so \(H\) is injective. A continuous bijection from the compact domain to the Hausdorff lift is a homeomorphism.

The entropy and measure assertions follow from \cref{thm:lyapunov}. Alternatively each fiber is parameterized by one radial interval with an equicontinuous family of iterates, giving a direct uniform finite Bowen cover. Under \(H\), the eventual core is precisely \(\Sigma\times\{0\}\); no assertion of surjectivity on \(\Sigma\times[0,1]\) is made.
\end{proof}

The full shift on \(M\geq2\) symbols is a continuous surjection of a compact metric space with positive topological entropy. The classical theorem \cite[Corollary 2.4(2)]{BGKM} applies on the apex subsystem and supplies Li--Yorke chaos there. Every one of its state projections is the same fixed state \(p\). In a common input fiber, even the full lift trajectories are asymptotic, because their input coordinates agree at all times and their state coordinates converge to \(p\). Thus total lift chaos and absence of state Li--Yorke pairs coexist without a contradiction.

\begin{center}
\begin{tabular}{@{}ll@{}}
\toprule
Quantity for the finite digit family & Value\\
\midrule
Exact invariance entropy on \(K_d\) or \(Q\) & \(\sum_i\log q_i\)\\
Outer and tracking invariance entropies & \(\sum_i\max\{0,\log(\rho q_i)\}\)\\
Total entropy of the true lift & \(\log M\)\\
Uniform entropy over the full input fibers & \(0\)\\
Entropy of the state-trajectory shift & \(0\)\\
Conditional KS entropy over input, any invariant measure & \(0\)\\
Entropy of any continuous feasible closed-loop state map & \(0\)\\
\bottomrule
\end{tabular}
\end{center}

\subsection{Feedback regularity and the control-set boundary}

\begin{proposition}[A feedback regularity obstruction]\label{prop:feedback}
For the cone system in \cref{thm:cone}, if \(\chi^+(A)>0\), no memoryless feedback \(\kappa:Q_B\to U\) continuous at \(p\) can satisfy \(f(x,\kappa(x))\in Q_B\) for every \(x\in Q_B\). If \(U\supset -AB\), there is nevertheless a Borel feedback whose induced state map is the restriction of a global smooth map,
\begin{align*}
 G(s,z)=(R(s),0).
\end{align*}
\end{proposition}
\begin{proof}
Suppose such a continuous-at-\(p\) feedback exists and write \(v_0=\kappa(p)\). For every \(y\in B\) and every \(0<s\leq1\), feasibility implies
\(Ay+\kappa(s,sy)\in B\).
Letting \(s\downarrow0\) and using closedness of \(B\) gives \(Ay+v_0\in B\) for all \(y\in B\). Thus the affine map \(T_0(y)=Ay+v_0\) sends \(B\) into itself. Iteration and subtraction give
\begin{align*}
 A^n(B-B)=T_0^n(B)-T_0^n(B)\subset B-B.
\end{align*}
Since \(B-B\) is bounded and contains a neighborhood of zero, the operators \(A^n\) are uniformly bounded. This excludes eigenvalues of modulus greater than one and contradicts \(\chi^+(A)>0\).

If \(U\supset -AB\), define \(\kappa(s,sy)=-Ay\) for \(s>0\) and \(\kappa(p)=0\). This is Borel and has values in \(U\). Substitution gives \(G(s,z)=(R(s),0)\) on \(Q_B\), which extends smoothly to the ambient space. Its feedback values depend on direction at the apex; continuity of the induced state map must not be confused with continuity of the feedback itself.
\end{proof}

For the nonlinear systems of \cref{thm:nonlinear}, the same obstruction to continuity of a viable feedback holds. Indeed feasibility at \((s,sy)\) would give
\begin{align*}
 qy+\kappa(s,sy)+P_{\kappa(s,sy)}(s,sy)/s\in[-1,1].
\end{align*}
The last term tends to zero uniformly in \(y\in[-1,1]\) and the finite input alphabet. Continuity at the apex would therefore force \(q[-1,1]+\kappa(p)\subset[-1,1]\), impossible for \(q>1\).

The common Lyapunov theorem excludes state Li--Yorke chaos for \emph{every} feasible controller in this family, not merely for the displayed resetting feedback. The scalar baseline in \cref{app:scalar} illustrates why a property of one chosen feedback is logically different from a conclusion quantified over all feasible inputs.

A control set ordinarily includes a controlled invariance condition together with approximate controllability: for every \(x\) in the set, the set is contained in the closure of the forward reachable set from \(x\), with a maximality condition. The cones here fail approximate controllability. The radial coordinate of every forward trajectory starting at \(s\in(0,1]\) is \(R^n(s)\leq s\), regardless of control and even if the other coordinates leave the cone. Points with larger radial coordinate cannot be approximated from such a state. Therefore \(Q_B\) is not a genuine control set. This is a substantive boundary, not a removable technical omission.

The equi-invariance/instability alternative discussed in the author version of \cite[Definition 3.3 and Theorem 3.4]{ZCH2021} does not repair the missing implication: its instability condition allows a nearby state under a fixed selected input to leave a neighborhood of the constraint. The two trajectories need not both remain feasible. In the cone family, strong expansion of an error in \(y=z/s\) is consistent with the fact that all \emph{feasible} ambient states eventually become close. Neither a single escape nor ordinary sensitivity supplies recurrent proximality.

Likewise, return to \(Q_B\) alone is not an additional recurrence hypothesis here, since feasible trajectories remain there at every time. A condition requiring repeated return to a nondegenerate interior region bounded away from the apex would exclude the present construction. Recurrence entropy \cite{SM2026} studies a different catalogue problem involving bounded-time returns; it cannot by itself be read as a statement about two feasible state trajectories.

The remaining focused question is whether positive invariance entropy on a genuine control set, together with explicitly stated accessibility and recurrence assumptions, forces a common-control or uniformly selected state scrambled structure. This paper proves neither such an implication nor its negation. A proof would need a source of recurrent proximality and persistent state separation on one fixed uncountable set. Merely obtaining a symbolic model with positive entropy, or allowing an arbitrary input shift, does not supply those two properties in the state projection.

\appendix
\section{Invariant partitions: what symbolic counting proves}\label{app:partition}

This appendix records the precise implication from an invariant partition and the two realizability failures relevant to the main question. It is supplementary material, not an additional novelty claim.

\begin{definition}
An invariant partition is a triple \(\mathcal C=(\mathcal P,\tau,w)\), where \(\mathcal P\) is a finite Borel partition of \(Q\), \(\tau\geq1\) is an integer, and \(w\) assigns to each cell \(P\) a control word \(w(P)\in U^\tau\) such that every \(x\in P\) stays in \(Q\) through times \(0,\ldots,\tau\) under that word. Define the sampled state map
\begin{align*}
 G_{\mathcal C}(x)=\varphi(\tau,x,w(P))\quad(x\in P).
\end{align*}
For \(x\in Q\), let \(I_{\mathcal C}(x)=(P_k)_{k\geq0}\), where \(G_{\mathcal C}^k(x)\in P_k\). Write
\(\Sigma_{\rm real}=I_{\mathcal C}(Q)\) and
\(\Sigma_{\mathcal C}=\overline{\Sigma_{\rm real}}\subset\mathcal P^{\mathbb N_0}\).
\end{definition}

\begin{proposition}\label{prop:partition}
The map \(G_{\mathcal C}\) is Borel, \(\Sigma_{\rm real}\) is forward shift invariant, and \(\Sigma_{\mathcal C}\) is a compact forward shift-invariant set. If \(N_n\) is the number of length-\(n\) prefixes in \(\Sigma_{\rm real}\), then
\begin{align*}
 r_{\mathrm{inv}}(n\tau,Q,Q)\leq N_n,\qquad
 h_{\mathrm{inv}}(Q,Q)\leq\tau^{-1}h_{\mathrm{top}}(\sigma|_{\Sigma_{\mathcal C}}).
\end{align*}
If the cells are clopen, \(G_{\mathcal C}\) and \(I_{\mathcal C}\) are continuous, \(\Sigma_{\rm real}=\Sigma_{\mathcal C}\), and
\begin{align*}
 \tau h_{\mathrm{inv}}(Q,Q)\leq
 h_{\mathrm{top}}(\sigma|_{\Sigma_{\mathcal C}})
 \leq h_{\mathrm{top}}(G_{\mathcal C}).
\end{align*}
\end{proposition}
\begin{proof}
On each Borel cell the formula for \(G_{\mathcal C}\) is a continuous finite-time solution map, so the finite pasting is Borel. The identity
\(I_{\mathcal C}(G_{\mathcal C}x)=\sigma I_{\mathcal C}(x)\)
gives forward invariance. Continuity of the shift passes that property to the closure.

For each realized prefix \((P_0,\ldots,P_{n-1})\), concatenate \(w(P_0),\ldots,w(P_{n-1})\) and extend arbitrarily to an infinite sequence. Every initial state having this itinerary prefix stays in \(Q\) through time \(n\tau\) under that one concatenated word. Thus \(N_n\) such controls span \(Q\). Taking closure creates no new finite prefixes: each prefix cylinder is open, so every cylinder meeting the closure already meets \(\Sigma_{\rm real}\).

The word counts are submultiplicative, \(N_{n+m}\leq N_nN_m\), and their normalized logarithms have a limit. The entropy of the one-sided subshift is this limit. Indeed distinct length-\(n\) prefixes give a uniformly separated set, while length-\((n+\ell)\) prefixes bound spanning numbers at a scale fixed by \(\ell\). Finally, for arbitrary horizon \(k\), monotonicity gives
\begin{align*}
 r_{\mathrm{inv}}(k,Q,Q)\leq
 r_{\mathrm{inv}}(\tau\lceil k/\tau\rceil,Q,Q)
 \leq N_{\lceil k/\tau\rceil}.
\end{align*}
This proves the sampling factor \(\tau^{-1}\) without discarding nonsampled horizons.

For clopen cells, finite pasting gives continuity of \(G_{\mathcal C}\). Each finite itinerary prefix has clopen preimage, so the itinerary map is continuous in the product topology. Its image is compact and hence closed; it is an onto factor map to \(\Sigma_{\mathcal C}\). The factor entropy inequality proves the last bound.
\end{proof}

If in the last case \(G_{\mathcal C}\) is also onto and \(h_{\mathrm{inv}}(Q,Q)>0\), the classical theorem \cite{BGKM} applies to \(G_{\mathcal C}\). This is a direct consequence of a known theorem under a continuous closed-loop hypothesis, not a new control-specific Li--Yorke principle. On connected \(Q\), a finite partition into nonempty clopen cells has just one cell. Its assigned word can be repeated forever from every state, so \(r_{\mathrm{inv}}(n,Q,Q)=1\) for all \(n\) and the exact entropy is zero. Thus that simple clopen route cannot apply to the positive-entropy solid cones.

\begin{remark}[A closure itinerary need not be realized]\label{rem:unrealized}
Let \(Q=[0,1]\), \(f(x,v)=2x+v\), and use the three cells
\begin{align*}
 P_0=[0,1/2),\quad P_1=[1/2,1),\quad P_2=\{1\},
\end{align*}
with one-step assigned controls \(0,-1,-2\), respectively. This is an invariant Borel partition. For each \(n\), the state \(x_n=1-2^{-(n+1)}\) remains in \(P_1\) for at least its first \(n\) sampled positions. Hence the constant sequence \(P_1P_1\cdots\) belongs to the itinerary closure. If it were realized by \(x\), iteration of \(x\mapsto2x-1\) would give
\(1-2^k(1-x)\in[1/2,1)\) for all \(k\), forcing \(x=1\). But \(1\in P_2\), a contradiction. Compactness of \(Q\) does not make a Borel itinerary map continuous.
\end{remark}

\begin{remark}[One-step graph edges need not concatenate]\label{rem:graph}
On the discrete compact space \(Q=\{a,b,c\}\), let a single control induce
\(a\mapsto b,\ b\mapsto b,\ c\mapsto a\).
For the partition \(A=\{a\}\), \(B=\{b,c\}\), the one-step intersection graph has edges \(A\to B\) and \(B\to A\). The graph therefore admits the path \(ABA\). Its only possible initial state is \(a\), whose actual orbit is \(a,b,b,\ldots\), so that path is unrealizable. Even an irreducible graph of one-step intersections can overcount actual paths. A Markov formula needs a separate path-realization hypothesis; it is not supplied by the existence of each edge. This is why the technical partition assumptions in \cite{ZHZ2023,YCYZ2025} cannot be omitted.
\end{remark}

\section{Scalar baselines and changes of control range}\label{app:scalar}

\subsection{Continuous time: an exact finite-horizon computation}

Only this subsection uses continuous time. Consider
\begin{align*}
 \dot x=\lambda x+u,\qquad \lambda>0,\quad
 u(t)\in[-\lambda,\lambda],\qquad Q=[-1,1].
\end{align*}
Controls are measurable essentially bounded functions. The variation-of-constants formula gives a unique solution for all \(t\geq0\). Define \(r_{\mathrm{inv}}(T,Q,Q)\) by the finite-catalogue quantifiers of \eqref{eq:quantifiers}, with every real \(t\in[0,T]\) in place of the integer indices; define exact and outer rates per unit of continuous time.

\begin{proposition}\label{prop:ct}
For every \(T\geq0\),
\begin{align*}
 r_{\mathrm{inv}}(T,Q,Q)=\lceil e^{\lambda T}\rceil,\qquad
 h_{\mathrm{inv}}(Q,Q)=h_{\mathrm{out}}(Q,Q)=\lambda.
\end{align*}
For any common input and distinct initial states the distance is
\(e^{\lambda t}|x-y|\), so no unconstrained common-input pair is proximal. The admissible feedback \(u=-\lambda x\) induces the identity flow on \(Q\).
\end{proposition}
\begin{proof}
For a fixed input, the time-\(T\) solution is an affine function of the initial state with slope \(e^{\lambda T}\). Initial states satisfying the final constraint therefore form an interval of length at most \(2e^{-\lambda T}\). Covering \(Q\), of length \(2\), requires at least \(\lceil e^{\lambda T}\rceil\) controls.

For \(T>0\), put \(\delta=e^{-\lambda T}\) and \(N=\lceil1/\delta\rceil\geq2\). Choose
\(c_j=-1+2j/(N-1)\), \(0\leq j<N\), and use the \(N\) constant inputs \(u_j=-\lambda c_j\). The associated solution is
\begin{align*}
 \varphi(t,x,u_j)=c_j+e^{\lambda t}(x-c_j).
\end{align*}
Its final state lies in \(Q\) exactly when
\begin{align*}
 x\in[(1-\delta)c_j-\delta,(1-\delta)c_j+\delta].
\end{align*}
These intervals cover \([-1,1]\): the first starts at \(-1\), the last ends at \(1\), and consecutive center gaps are
\(2(1-\delta)/(N-1)\leq2\delta\).
For each covered initial state the trajectory lies between its initial and final values, so remains in \(Q\) for the whole real interval \([0,T]\). This proves the upper bound. At \(T=0\), one arbitrary input suffices.

For \(Q^\varepsilon=(-1-\varepsilon,1+\varepsilon)\), the final-time lower bound becomes
\begin{align*}
 r_{\mathrm{inv}}(T,Q,Q^\varepsilon)\geq
             e^{\lambda T}/(1+\varepsilon).
\end{align*}
The exact catalogue remains an outer catalogue, proving outer rate \(\lambda\) for each fixed \(\varepsilon>0\). Subtracting two solutions with the same input proves their distance formula. Finally \(|-\lambda x|\leq\lambda\) on \(Q\), and the stated feedback gives \(\dot x=0\).
\end{proof}

Two distinct initial states cannot share an infinitely feasible input in this bounded interval, since their distance would become unbounded. The stronger unconstrained calculation in \cref{prop:ct} prevents the common-control conclusion from being a merely vacuous consequence of feasibility.
More generally, an autonomous Lipschitz scalar closed-loop vector field preserving a compact interval has convergent trajectories. A nonstationary solution cannot cross an equilibrium by uniqueness, so its derivative has a constant sign along its orbit and it is bounded and monotone. Pair distances therefore have limits, precluding Li--Yorke pairs. This observation is about autonomous scalar feedbacks; it is not a theorem about all time-dependent controllers.

\subsection{Interval controls and interior containment in the cone}

In the planar digit cone, enlarge \(U_q\) to its convex hull
\(\widehat U_q=[-(q-1),q-1]\).
The exact lower bounds in \cref{thm:digits} use only the slope of a fixed-word map, and remain valid for this enlarged range. The original digit controls still give the upper bounds. Thus all the exact rates and counts in the planar case remain unchanged. \Cref{thm:cone}(ii) likewise gives the same outer and tracking rates.

For initial states in \(K_1\), the exact rate is unchanged if strict interior containment is required at every positive integer time. Indeed their ratios \(y=z/s\) lie in \([-1/2,1/2]\). Cover this interval by finitely many open intervals of radius \(q^{-n}\), with centers \(c\in[-1/2,1/2]\), using \(O(q^n)\) centers. For each center take \(u_0=-qc\), \(u_k=0\) for \(k\geq1\). These controls lie in \(\widehat U_q\), since \(q/2\leq q-1\). For an assigned initial ratio,
\begin{align*}
 |y_k|=q^k|y-c|<1\quad(1\leq k\leq n).
\end{align*}
The radial coordinate remains strictly between \(0\) and \(1\), so the ambient states lie in \(\operatorname{int} Q\). This proves an upper rate \(\log q\), while exact containment gives the matching lower bound. This assertion is for \(K_1\); requiring interior containment from a boundary point of the whole cone is a different question.

Finally, the continuous state map \((s,z)\mapsto(R(s),t(s)z)\) is induced on \(Q\) by
\(\kappa(s,z)=-(q-1)z/s\) for \(s>0\), with \(\kappa(p)=0\).
It is feasible and smooth as a state map, although \(\kappa\) is discontinuous at \(p\). This gives a zero-entropy continuous closed loop using the smaller convexified digit range, without the larger range needed for one-step resetting.

\begin{thebibliography}{99}
\small\setlength{\itemsep}{4pt}

\bibitem{BGKM}
F. Blanchard, E. Glasner, S. Kolyada, and A. Maass,
On Li-Yorke pairs, \emph{J. Reine Angew. Math.} \textbf{547} (2002), 51--68.
\href{https://doi.org/10.1515/crll.2002.053}{doi:\nolinkurl{10.1515/crll.2002.053}}.

\bibitem{CK2009}
F. Colonius and C. Kawan,
Invariance entropy for control systems, \emph{SIAM J. Control Optim.} \textbf{48} (2009), no.~3, 1701--1721.
\href{https://doi.org/10.1137/080713902}{doi:\nolinkurl{10.1137/080713902}}.

\bibitem{K2009}
C. Kawan, \emph{Invariance Entropy for Control Systems},
doctoral dissertation, Universit\"at Augsburg, 2009.
\href{https://d-nb.info/1000462749/34}{German National Library full text}.

\bibitem{K2011}
C. Kawan, Invariance entropy of control sets,
\emph{SIAM J. Control Optim.} \textbf{49} (2011), no.~2, 732--751.
\href{https://doi.org/10.1137/100783340}{doi:\nolinkurl{10.1137/100783340}}.

\bibitem{KStrict2011}
C. Kawan, Lower bounds for the strict invariance entropy,
\emph{Nonlinearity} \textbf{24} (2011), no.~7, 1909--1935.
\href{https://doi.org/10.1088/0951-7715/24/7/001}{doi:\nolinkurl{10.1088/0951-7715/24/7/001}}.

\bibitem{DSK2016}
A. Da Silva and C. Kawan, Invariance entropy of hyperbolic control sets,
\emph{Discrete Contin. Dyn. Syst.} \textbf{36} (2016), no.~1, 97--136.
\href{https://doi.org/10.3934/dcds.2016.36.97}{doi:\nolinkurl{10.3934/dcds.2016.36.97}}. Author version: \href{https://arxiv.org/abs/1408.2416v1}{arXiv:\nolinkurl{1408.2416v1}}.

\bibitem{KDS2018}
C. Kawan and A. Da Silva, Invariance entropy for a class of partially hyperbolic sets,
\emph{Math. Control Signals Systems} \textbf{30} (2018), no.~4, article 18.
\href{https://doi.org/10.1007/s00498-018-0224-2}{doi:\nolinkurl{10.1007/s00498-018-0224-2}}. Author version: \href{https://arxiv.org/abs/1711.01181v2}{arXiv:\nolinkurl{1711.01181v2}}.

\bibitem{KChaos}
C. Kawan, Control of chaos with minimal information transfer,
preprint, version 4, 2021. \href{https://arxiv.org/abs/2003.06935v4}{arXiv:\nolinkurl{2003.06935v4}}.

\bibitem{Yao2026}
S. Yao, Invariance entropy in the dust,
preprint, version 1, 2026. \href{https://arxiv.org/abs/2607.02279v1}{arXiv:\nolinkurl{2607.02279v1}}.

\bibitem{CH2021}
F. Colonius and B. Hamzi, Entropy for practical stabilization,
\emph{SIAM J. Control Optim.} \textbf{59} (2021), no.~3, 2195--2222.
\href{https://doi.org/10.1137/20M1367775}{doi:\nolinkurl{10.1137/20M1367775}}. Author version: \href{https://arxiv.org/abs/2009.08187v3}{arXiv:\nolinkurl{2009.08187v3}}.

\bibitem{KD2016}
C. Kawan and J.-C. Delvenne, Network entropy and data rates required for networked control,
\emph{IEEE Trans. Control Netw. Syst.} \textbf{3} (2016), no.~1, 57--66.
\href{https://doi.org/10.1109/TCNS.2015.2440551}{doi:\nolinkurl{10.1109/TCNS.2015.2440551}}. Author version: \href{https://arxiv.org/abs/1409.6037v1}{arXiv:\nolinkurl{1409.6037v1}}.

\bibitem{CST2023}
F. Colonius, A. J. Santana, and E. C. Viscovini,
Chain controllability of linear control systems,
\emph{SIAM J. Control Optim.} \textbf{62} (2024), no.~4, 2387--2411.
\href{https://doi.org/10.1137/23M1626347}{doi:\nolinkurl{10.1137/23M1626347}}. Author version: \href{https://arxiv.org/abs/2312.14595v1}{arXiv:\nolinkurl{2312.14595v1}}.

\bibitem{CCS2021}
F. Colonius, J. A. N. Cossich, and A. J. Santana,
Outer invariance entropy for discrete-time linear systems on Lie groups,
\emph{ESAIM Control Optim. Calc. Var.} \textbf{27} (2021), article 55.
\href{https://doi.org/10.1051/cocv/2021054}{doi:\nolinkurl{10.1051/cocv/2021054}}. Author version: \href{https://arxiv.org/abs/2008.03219v2}{arXiv:\nolinkurl{2008.03219v2}}.

\bibitem{ZCH2021}
X.-F. Zhong, Z.-J. Chen, and Y. Huang,
Equi-invariability, bounded invariance complexity and L-stability for control systems,
\emph{Sci. China Math.} \textbf{64} (2021), no.~10, 2275--2294.
\href{https://doi.org/10.1007/s11425-020-1693-7}{doi:\nolinkurl{10.1007/s11425-020-1693-7}}. Author version: \href{https://arxiv.org/abs/2005.10457v1}{arXiv:\nolinkurl{2005.10457v1}}.

\bibitem{ZC2023}
X.-F. Zhong and Z.-J. Chen, Equi-invariability and bounded invariance complexity for control systems,
\emph{J. Dyn. Differential Equations} \textbf{35} (2023), no.~3, 2261--2277.
\href{https://doi.org/10.1007/s10884-022-10161-2}{doi:\nolinkurl{10.1007/s10884-022-10161-2}}.

\bibitem{CZ2024}
Z. Chen and X. Zhong, Invariance complexity and equi-invariability for control systems,
\emph{J. Math. Anal. Appl.} \textbf{539} (2024), no.~2, article 128533.
\href{https://doi.org/10.1016/j.jmaa.2024.128533}{doi:\nolinkurl{10.1016/j.jmaa.2024.128533}}.

\bibitem{WH2024}
T. Wang and Y. Huang, The measurable data rate theorem for discrete-time control systems,
\emph{J. Dyn. Differential Equations} \textbf{36} (2024), no.~4, 3259--3279.
\href{https://doi.org/10.1007/s10884-023-10269-z}{doi:\nolinkurl{10.1007/s10884-023-10269-z}}.

\bibitem{ZH2025}
X. Zhong and Y. Huang, Dimension types of invariance entropies for uncertain control systems,
\emph{Sci. China Math.} \textbf{68} (2025), no.~12, 2917--2932.
\href{https://doi.org/10.1007/s11425-024-2422-2}{doi:\nolinkurl{10.1007/s11425-024-2422-2}}.

\bibitem{CHZ2026}
H. Chen, Y. Huang, and X. Zhong, Variational principles of invariance entropy dimension,
\emph{J. Differential Equations} \textbf{453} (2026), article 113819.
\href{https://doi.org/10.1016/j.jde.2025.113819}{doi:\nolinkurl{10.1016/j.jde.2025.113819}}.

\bibitem{Zhang2006}
G. Zhang, Relative entropy, asymptotic pairs and chaos,
\emph{J. London Math. Soc.} \textbf{73} (2006), no.~1, 157--172.
\href{https://doi.org/10.1112/S0024610705022520}{doi:\nolinkurl{10.1112/S0024610705022520}}.

\bibitem{LTZ2024}
C. Liu, F. Tan, and J. Zhang,
Mean Li-Yorke chaos along any infinite sequence for infinite-dimensional random dynamical systems,
\emph{J. Differential Equations} \textbf{403} (2024), 548--575.
\href{https://doi.org/10.1016/j.jde.2024.05.035}{doi:\nolinkurl{10.1016/j.jde.2024.05.035}}. Author version: \href{https://arxiv.org/abs/2207.08505v3}{arXiv:\nolinkurl{2207.08505v3}}.

\bibitem{YCYZ2025}
R. Yang, E. Chen, J. Yang, and X. Zhou,
Bowen's equations for invariance pressure of control systems,
\emph{SIAM J. Control Optim.} \textbf{63} (2025), no.~2, 1104--1128.
\href{https://doi.org/10.1137/23M1607684}{doi:\nolinkurl{10.1137/23M1607684}}. Author version: \href{https://arxiv.org/abs/2309.01628v3}{arXiv:\nolinkurl{2309.01628v3}}.

\bibitem{ZHZ2023}
X. Zhong, Y. Huang, and X. Zou,
Invariance entropy for uncertain control systems,
\emph{Systems Control Lett.} \textbf{180} (2023), article 105623.
\href{https://doi.org/10.1016/j.sysconle.2023.105623}{doi:\nolinkurl{10.1016/j.sysconle.2023.105623}}.

\bibitem{SM2026}
H. Sibai and E. Mallada,
Recurrence of nonlinear control systems: Entropy, bit rates, and finite alphabet controllers,
\emph{Nonlinear Anal. Hybrid Syst.} \textbf{59} (2026), article 101649.
\href{https://doi.org/10.1016/j.nahs.2025.101649}{doi:\nolinkurl{10.1016/j.nahs.2025.101649}}.

\bibitem{CF2026}
F. Colonius and R. Fabbri,
Nonautonomous control systems and skew product flows,
\emph{J. Dyn. Control Syst.} \textbf{32} (2026), no.~1, article 3.
\href{https://doi.org/10.1007/s10883-025-09750-3}{doi:\nolinkurl{10.1007/s10883-025-09750-3}}.

\bibitem{WHS2019}
T. Wang, Y. Huang, and H.-W. Sun,
Measure-theoretic invariance entropy for control systems,
\emph{SIAM J. Control Optim.} \textbf{57} (2019), no.~1, 310--333.
\href{https://doi.org/10.1137/18M1197862}{doi:\nolinkurl{10.1137/18M1197862}}.

\bibitem{NWH2022}
X. Nie, T. Wang, and Y. Huang,
Measure-theoretic invariance entropy and variational principles for control systems,
\emph{J. Differential Equations} \textbf{321} (2022), 318--348.
\href{https://doi.org/10.1016/j.jde.2022.03.013}{doi:\nolinkurl{10.1016/j.jde.2022.03.013}}.

\bibitem{CQuasi2018}
F. Colonius, Invariance entropy, quasi-stationary measures and control sets,
\emph{Discrete Contin. Dyn. Syst.} \textbf{38} (2018), no.~4, 2093--2123.
\href{https://doi.org/10.3934/dcds.2018086}{doi:\nolinkurl{10.3934/dcds.2018086}}. Author version: \href{https://arxiv.org/abs/1705.08658}{arXiv:\nolinkurl{1705.08658}}.
\end{thebibliography}
\end{document}
